% Options for packages loaded elsewhere
\PassOptionsToPackage{unicode}{hyperref}
\PassOptionsToPackage{hyphens}{url}
\PassOptionsToPackage{dvipsnames,svgnames,x11names}{xcolor}
%
\documentclass[
  10pt,
]{article}
\usepackage{amsmath,amssymb}
\usepackage{iftex}
\ifPDFTeX
  \usepackage[T1]{fontenc}
  \usepackage[utf8]{inputenc}
  \usepackage{textcomp} % provide euro and other symbols
\else % if luatex or xetex
  \usepackage{unicode-math} % this also loads fontspec
  \defaultfontfeatures{Scale=MatchLowercase}
  \defaultfontfeatures[\rmfamily]{Ligatures=TeX,Scale=1}
\fi
\usepackage{lmodern}
\ifPDFTeX\else
  % xetex/luatex font selection
\fi
% Use upquote if available, for straight quotes in verbatim environments
\IfFileExists{upquote.sty}{\usepackage{upquote}}{}
\IfFileExists{microtype.sty}{% use microtype if available
  \usepackage[]{microtype}
  \UseMicrotypeSet[protrusion]{basicmath} % disable protrusion for tt fonts
}{}
\makeatletter
\@ifundefined{KOMAClassName}{% if non-KOMA class
  \IfFileExists{parskip.sty}{%
    \usepackage{parskip}
  }{% else
    \setlength{\parindent}{0pt}
    \setlength{\parskip}{6pt plus 2pt minus 1pt}}
}{% if KOMA class
  \KOMAoptions{parskip=half}}
\makeatother
\usepackage{xcolor}
\usepackage{longtable,booktabs,array}
\usepackage{calc} % for calculating minipage widths
% Correct order of tables after \paragraph or \subparagraph
\usepackage{etoolbox}
\makeatletter
\patchcmd\longtable{\par}{\if@noskipsec\mbox{}\fi\par}{}{}
\makeatother
% Allow footnotes in longtable head/foot
\IfFileExists{footnotehyper.sty}{\usepackage{footnotehyper}}{\usepackage{footnote}}
\makesavenoteenv{longtable}
\setlength{\emergencystretch}{3em} % prevent overfull lines
\providecommand{\tightlist}{%
  \setlength{\itemsep}{0pt}\setlength{\parskip}{0pt}}
\setcounter{secnumdepth}{5}

\usepackage{amsmath,amssymb}
\usepackage{graphicx}
\usepackage{longtable,booktabs,array}
\usepackage{xcolor}
\usepackage[margin=2.0cm]{geometry}
\usepackage{microtype}
\setlength{\emergencystretch}{3em}
\setcounter{MaxMatrixCols}{20}
\pdfinfoomitdate=1
\pdftrailerid{}
\pdfsuppressptexinfo=-1
\renewcommand{\arraystretch}{1.12}
\usepackage{fvextra}
\RecustomVerbatimEnvironment{verbatim}{Verbatim}{breaklines,breakanywhere,fontsize=\small}
\fvset{breaklines,breakanywhere}
\allowdisplaybreaks
\sloppy
\ifLuaTeX
  \usepackage{selnolig}  % disable illegal ligatures
\fi
\IfFileExists{bookmark.sty}{\usepackage{bookmark}}{\usepackage{hyperref}}
\IfFileExists{xurl.sty}{\usepackage{xurl}}{} % add URL line breaks if available
\urlstyle{same}
\hypersetup{
  pdftitle={dirac matrices: the eight real 16 x 16 Dirac matrices of the author's notebook 4+4.nb},
  pdfauthor={Matrices and notebook: Patrick L. Nash; Provenance: Claude (Anthropic), 2026-10-24},
  colorlinks=true,
  linkcolor={blue},
  filecolor={Maroon},
  citecolor={Blue},
  urlcolor={blue},
  pdfcreator={LaTeX via pandoc}}

\title{dirac matrices: the eight real 16 x 16 Dirac matrices of the
author's notebook 4+4.nb}
\usepackage{etoolbox}
\makeatletter
\providecommand{\subtitle}[1]{% add subtitle to \maketitle
  \apptocmd{\@title}{\par {\large #1 \par}}{}{}
}
\makeatother
\subtitle{Provenance file of the model UM-4 (universe\_model\_4):
construction, reality, anticommutation relations, spin(4,4), Spin(4,4)
and Pin(4,4), the chiral projection operators, and the proof that every
calculation uses these matrices}
\author{Matrices and notebook: Patrick L. Nash \and Provenance: Claude
(Anthropic), 2026-10-24}
\date{}

\begin{document}
\maketitle

{
\hypersetup{linkcolor=}
\setcounter{tocdepth}{3}
\tableofcontents
}
\hypertarget{the-answer}{%
\section{The answer}\label{the-answer}}

This file was generated by the program
\texttt{um4\_dirac\_matrices\_doc.py} of UM-4 from the results of
\texttt{um4\_dirac\_matrices.py}. It proves, with exact computations
whose results are printed here, the statements the author asked for on
2026-10-24 about the eight real \(16 \times 16\) Dirac matrices. It is
self-contained: every equation it uses is written here and every symbol
is defined right below the equation in which it appears. All 55 named
checks of the program
\texttt{universe\_model\_4/code/um4\_dirac\_matrices.py} pass; each
statement below names the check that proves it.

\begin{enumerate}
\def\labelenumi{\arabic{enumi}.}
\tightlist
\item
  \textbf{They are the matrices of \texttt{4+4.nb}.} The eight matrices
  \(T_a\) = \texttt{T16\^{}A{[}a{]}} (\(a = 0, \dots, 7\)) were rebuilt
  cell by cell from the author's own input cells (quoted verbatim in
  Section 2) and are equal, entry by entry, to every matrix that the
  notebook itself stored as an output when Mathematica evaluated those
  cells: the eight \(T_a\) (\texttt{Out{[}210{]}}),
  \(T_0 T_1 \cdots T_7\) (\texttt{Out{[}98{]}}), \(P_L\) and \(P_R\)
  (\texttt{Out{[}111{]}}, \texttt{Out{[}112{]}}) and seven more; 26
  stored True/False outputs (283 values) of the matrix cells are
  reproduced identically. \texttt{B11\_stored\_T16A\_0\_to\_7} (passed),
  \texttt{B8\_stored\_T16A8} (passed),
  \texttt{B13\_stored\_boolean\_outputs} (passed).
\item
  \textbf{They are the matrices used in this repository's calculations.}
  The same eight matrices, entry by entry, are used by the author's
  notebook
  \texttt{Pair\_Creation\_of\_Universes\_WaveFunctionOfUniverse-4+4-Einstein-Lovelock-Nash.nb}
  (the same defining formulas), the Revision record, the universe model
  UM-2, nine Wolfram packages and the Stage-1 fixture (Section 11);
  every UM-4 computation takes them from one module,
  \texttt{universe\_model\_4/code/um4\_gammas.py}.
  \texttt{D1\_pair\_creation\_extraction} (passed),
  \texttt{D2\_pair\_creation\_same\_formulas} (passed),
  \texttt{D5\_repository\_wolfram\_packages} (passed).
\item
  \textbf{They are real.} Every entry is \(-1\), \(0\) or \(+1\); each
  \(T_a\) is a signed permutation matrix and an orthogonal matrix,
  symmetric for \(a = 0, 1, 2, 3\) and antisymmetric for
  \(a = 4, 5, 6, 7\). \texttt{C1\_real\_integer\_entries} (passed),
  \texttt{C2\_signed\_permutation} (passed),
  \texttt{C4\_orthogonal\_and\_symmetry} (passed).
\item
  \textbf{They satisfy the anticommutation relations}
  \(\{T_a, T_b\} = 2\eta_{ab}I_{16}\) with
  \(\eta = \mathrm{diag}(+1, +1, +1, +1, -1, -1, -1, -1)\): they
  generate the real Clifford algebra of signature \((4,4)\), which is
  the full algebra of real \(16 \times 16\) matrices.
  \texttt{C3\_clifford\_relations} (passed), \texttt{C30\_basis\_256}
  (passed).
\item
  \textbf{Their scaled commutators generate Pin(4,4).} The 28 matrices
  \(S_{ab} = \tfrac14[T_a, T_b]\) (\(a < b\)) are linearly independent,
  close under commutation with the structure constants of \(so(4,4)\)
  and act on the \(T_c\) as the vector representation: they are a basis
  of the Lie algebra \(spin(4,4) \cong so(4,4)\), which is the Lie
  algebra of Pin(4,4) (and of Spin(4,4)). Each exponential
  \(\exp(\theta S_{ab})\) is a product of two unit vectors, so it lies
  in Pin(4,4); the exponentials generate the identity component
  \(\mathrm{Spin}_0(4,4)\); together with the two reflections \(T_0\)
  and \(T_4\) they generate all of Pin(4,4), whose four components are
  identified by exact invariants. Precisely stated: the scaled
  commutators generate Pin(4,4) as its Lie algebra and, by
  exponentiation, its identity component; no product of exponentials
  leaves that component, so the reflections are needed for the other
  three. \texttt{C21\_SAB\_independent} (passed),
  \texttt{C24\_spin44\_closure} (passed),
  \texttt{C26\_exponentials\_two\_unit\_vectors} (passed),
  \texttt{C28\_exponentials\_in\_identity\_component} (passed),
  \texttt{C27\_pin44\_components} (passed).
\item
  \textbf{The chiral projection operators of the subsection ``Left -
  Handed vs . Right - Handed Fermion Fields (Chirality) ; projection
  operators'' are re-derived and confirmed.} The only matrices that
  commute with every \(S_{ab}\) are the combinations of \(I_{16}\) and
  \(\Gamma = T_0 T_1 \cdots T_7\); \(\Gamma^2 = I_{16}\); the only
  non-trivial projectors in that algebra are
  \(\tfrac12(I_{16} \mp \Gamma)\), which are the author's \(P_L\) and
  \(P_R\), entry by entry; \(\Gamma = \mathrm{diag}(-I_8, I_8)\), so
  \(P_L\) projects on \(\Psi_0, \dots, \Psi_7\) and \(P_R\) on
  \(\Psi_8, \dots, \Psi_{15}\). \texttt{C17\_commutant\_spin44}
  (passed), \texttt{C20\_idempotents\_of\_chirality\_algebra} (passed),
  \texttt{B9\_stored\_PL} (passed), \texttt{B10\_stored\_PR} (passed).
\item
  \textbf{The coordinates are the author's}
  \(\{x_0, x_1, x_2, x_3, x_4, x_5, x_6, x_7\}\): the input cell
  \texttt{In{[}48{]}} of the Section DEFCHART of
  \texttt{Dirac\_Einstein-Lovelock\_4+4\_Nash.nb} is
  \texttt{X=\{x0{[}{]},\ x1{[}{]},\ x2{[}{]},\ x3{[}{]},\ x4{[}{]},\ x5{[}{]},\ x6{[}{]},\ x7{[}{]}\}}
  (spaces added after the commas), and \texttt{4+4.nb} uses the same
  list (its cell 9). The frame index \(a\) of \(T_a\) is the index of
  the coordinate \(x_a\). \texttt{A6\_coordinates\_DEFCHART} (passed),
  \texttt{A5\_coordinates\_in\_4p4} (passed).
\end{enumerate}

Each of the eight matrices, the chirality matrix, the Dirac-adjoint
matrix, the projectors, the complex structure and the 28 generators is
displayed in full below.

\hypertarget{notation}{%
\section{Notation}\label{notation}}

The frame metric of the eight-dimensional space is

\[
\eta_{ab} = \mathrm{diag}(+1,\ +1,\ +1,\ +1,\ -1,\ -1,\ -1,\ -1)_{ab},\qquad a, b = 0, 1, \dots, 7
\]

\textbf{Symbols in this equation.}

\begin{itemize}
\tightlist
\item
  \(\eta_{ab}\): The flat frame metric,
  \(\eta = \mathrm{diag}(+1, +1, +1, +1, -1, -1, -1, -1)\) in the order
  \(x_0, \dots, x_7\) (the author's \texttt{eta4488}); \(\eta^{ab}\) is
  the same matrix.
\item
  \(a, b, c, d\): Frame (orthonormal) indices running over
  \(0, 1, \dots, 7\); the value \(k\) belongs to the author's coordinate
  \(x_k\) (\(x_0\) the hidden space-like direction, \(x_1, x_2, x_3\)
  ordinary space, \(x_4\) the time, \(x_5, x_6, x_7\) the extra
  time-like directions). No sum over a repeated frame index unless a sum
  is written.
\end{itemize}

Four directions are space-like (\(x_0\), \(x_1\), \(x_2\), \(x_3\)) and
four are time-like (\(x_4\), \(x_5\), \(x_6\), \(x_7\)), in the order of
the author's coordinates:

\[
X = \{x_0,\ x_1,\ x_2,\ x_3,\ x_4,\ x_5,\ x_6,\ x_7\}
\]

\textbf{Symbols in this equation.}

\begin{itemize}
\tightlist
\item
  \(x_0, \dots, x_7\): The author's coordinates
  \(X = \{x_0, x_1, \dots, x_7\}\) (notebook
  \texttt{Dirac\_Einstein-Lovelock\_4+4\_Nash.nb}, Section DEFCHART,
  \texttt{X\ =\ \{x0{[}{]},\ x1{[}{]},\ ...,\ x7{[}{]}\}}).
\end{itemize}

Matrices act on columns
\(\Psi = (\Psi_0, \dots, \Psi_{15})^{\mathsf T}\); rows and columns of a
\(16 \times 16\) matrix are counted from 0 in the text and in the
displays (the notebook counts them from 1). A block matrix
\(\begin{pmatrix} A & B \\ C & D \end{pmatrix}\) is the author's
\texttt{ArrayFlatten{[}\{\{A,\ B\},\ \{C,\ D\}\}{]}}.

\hypertarget{the-source-the-authors-notebook-44.nb}{%
\section{The source: the author's notebook
4+4.nb}\label{the-source-the-authors-notebook-44.nb}}

The notebook \texttt{universe\_model\_3/author\_inputs/4+4.nb} (``eight,
real-valued, 16 by 16, Dirac Matrices'', Wolfram 15.0) has the sha256
\texttt{7ec1709a0028bce0b443f116cf27b5d3e670de7aa195659a9c3274527270429d}
and 806 cells; it is read without Mathematica by
\texttt{universe\_model\_4/tools/nb\_reader.py}, which parses the
notebook file (a Wolfram Language expression) and converts every cell to
linear text. \texttt{A1\_notebook\_sha256} (passed),
\texttt{A2\_notebook\_cell\_count} (passed).

The matrices are defined by the following input cells, quoted verbatim
(cell number in the file, label, text;
\texttt{A3\_defining\_cells\_verbatim} (passed)):

\begin{verbatim}
cell 10  In[7]:=
    ID4=IdentityMatrix[4];
    ID8=IdentityMatrix[8];
cell 11  In[9]:=
    (\[Eta]4488 = ArrayFlatten[{{IdentityMatrix[4], 0}, {0, -IdentityMatrix[4]}}]) // MatrixForm
cell 16  In[18]:=
    \[Sigma]=\[Sigma]88=ArrayFlatten[{{ConstantArray[0,{4,4}],IdentityMatrix[4]},{IdentityMatrix[4],ConstantArray[0,{4,4}]}}];Protect[\[Sigma]88];
cell 17  In[19]:=
    Zero16=ConstantArray[0,{16,16}];
    G16=ArrayFlatten[{{\[Eta]8,0},{0,-\[Eta]8}}];
    ID16=Id16=IdentityMatrix[16];MId16=ExpandAll[-1*Id16];
    Simp16={Zero16->ZERO16,Id16->ONE16,MId16->MONE16};
cell 24  In[29]:=
    Qa[h_,p_,q_]:=Signature[{h,p,q,4}];
    Qb[h_,p_,q_]:= ID4[[p,4]]*ID4[[q,h]]-ID4[[p,h]]*ID4[[q,4]];
    SelfDualAntiSymmetric[h_,p_,q_]:=Qa[h,p,q] -Qb[h,p,q] ;
    AntiSelfDualAntiSymmetric[h_,p_,q_]:= (Qa[h,p,q] + Qb[h,p,q] );
cell 29  In[35]:=
    Do[s4by4[h]=Table[Table[SelfDualAntiSymmetric[h,p,q], {q, 4}],{p, 4}], {h,1, 3} ]
cell 30  In[36]:=
    Do[t4by4[h]=Table[Table[AntiSelfDualAntiSymmetric[h,p,q], {q, 4}],{p, 4}], {h,1, 3} ]
cell 73  In[60]:=
    \[Zeta][0]=ID8;
    Table[\[Zeta][7-h]= ArrayFlatten[{{0,t4by4[h]},{-t4by4[h],0}}], {h,1, 3} ];
    Table[\[Zeta][h]= ArrayFlatten[{{0,s4by4[h]},{s4by4[h],0}}], {h,1, 3} ];
    (\[Zeta][7]=\[Zeta][1].\[Zeta][2].\[Zeta][3].\[Zeta][4].\[Zeta][5].\[Zeta][6])//MatrixForm
cell 87  In[70]:=
    \[Tau][0]=ID8;
    Do[{\[Tau][A] = FullSimplify[ExpandAll[(\[Zeta][A])]]} , {A, 1, 7}]
    (*Do[{\[Tau][A] = FullSimplify[ExpandAll[-(\[Sigma].\[Zeta][1].\[Zeta][2].\[Zeta][3].\[Zeta][A])]]} , {A, 1, 7}]*)
cell 91  In[74]:=
    Overscript[\[Tau], _][0] =ID8;
    Do[{Overscript[\[Tau], _][A] = FullSimplify[ExpandAll[(\[Sigma].Transpose[\[Tau][A]].\[Sigma])]]} , {A, 1, 7}]
cell 131  In[96]:=
    (\[Sigma]16=T16^A[0].T16^A[1].T16^A[2].T16^A[3])//MatrixForm
cell 133  In[97]:=
    Table[T16^A[A1]= ArrayFlatten[{{0,Overscript[\[Tau], _][A1]},{\[Tau][A1],0}}], {A1,0, 7} ];
cell 134  In[98]:=
    (T16^A[8]=T16^A[0].T16^A[1].T16^A[2].T16^A[3].T16^A[4].T16^A[5].T16^A[6].T16^A[7])//MatrixForm
cell 161  In[111]:=
    (Subscript[P, L]=1/2(ID16-T16^A[8]))//MatrixForm
cell 163  In[112]:=
    (Subscript[P, R]=1/2(ID16+T16^A[8]))//MatrixForm
\end{verbatim}

The cells 160 to 166 form the subsection ``Left - Handed vs . Right -
Handed Fermion Fields (Chirality) ; projection operators''
(\texttt{A4\_chirality\_subsection\_found} (passed)): cell 161 defines
\(P_L\), cell 163 defines \(P_R\) and cell 165 tests them; Section 9
below re-derives them.

In cell 131 (\texttt{In{[}96{]}}) the author assigns
\texttt{sigma16\ =\ T16\^{}A{[}0{]}.T16\^{}A{[}1{]}.T16\^{}A{[}2{]}.T16\^{}A{[}3{]}}
before the matrices are defined in cell 133 (\texttt{In{[}97{]}}): the
symbol then holds the unevaluated product, which takes its value as soon
as the \texttt{T16\^{}A} exist; cell 148 (\texttt{In{[}105{]}}) confirms
the value \(\mathrm{diag}(-\sigma, \sigma)\). In cell 87
(\texttt{In{[}70{]}}) the line inside \texttt{(*\ ...\ *)} is a comment
and is not evaluated.

\hypertarget{the-construction-step-by-step}{%
\section{The construction, step by
step}\label{the-construction-step-by-step}}

\hypertarget{the-two-8-times-8-metrics}{%
\subsection{\texorpdfstring{The two \(8 \times 8\)
metrics}{The two 8 \textbackslash times 8 metrics}}\label{the-two-8-times-8-metrics}}

\[
\eta = \begin{pmatrix} I_4 & 0 \\ 0 & -I_4 \end{pmatrix} = \begin{pmatrix} 1 & 0 & 0 & 0 & 0 & 0 & 0 & 0 \\ 0 & 1 & 0 & 0 & 0 & 0 & 0 & 0 \\ 0 & 0 & 1 & 0 & 0 & 0 & 0 & 0 \\ 0 & 0 & 0 & 1 & 0 & 0 & 0 & 0 \\ 0 & 0 & 0 & 0 & -1 & 0 & 0 & 0 \\ 0 & 0 & 0 & 0 & 0 & -1 & 0 & 0 \\ 0 & 0 & 0 & 0 & 0 & 0 & -1 & 0 \\ 0 & 0 & 0 & 0 & 0 & 0 & 0 & -1 \end{pmatrix},\qquad \sigma = \begin{pmatrix} 0 & I_4 \\ I_4 & 0 \end{pmatrix} = \begin{pmatrix} 0 & 0 & 0 & 0 & 1 & 0 & 0 & 0 \\ 0 & 0 & 0 & 0 & 0 & 1 & 0 & 0 \\ 0 & 0 & 0 & 0 & 0 & 0 & 1 & 0 \\ 0 & 0 & 0 & 0 & 0 & 0 & 0 & 1 \\ 1 & 0 & 0 & 0 & 0 & 0 & 0 & 0 \\ 0 & 1 & 0 & 0 & 0 & 0 & 0 & 0 \\ 0 & 0 & 1 & 0 & 0 & 0 & 0 & 0 \\ 0 & 0 & 0 & 1 & 0 & 0 & 0 & 0 \end{pmatrix}
\]

\textbf{Symbols in this equation.}

\begin{itemize}
\tightlist
\item
  \(\eta_{ab}\): The flat frame metric,
  \(\eta = \mathrm{diag}(+1, +1, +1, +1, -1, -1, -1, -1)\) in the order
  \(x_0, \dots, x_7\) (the author's \texttt{eta4488}); \(\eta^{ab}\) is
  the same matrix.
\item
  \(\sigma\): The real symmetric \(8 \times 8\) matrix
  \(\sigma = \begin{pmatrix} 0 & I_4 \\ I_4 & 0 \end{pmatrix}\) (the
  author's \texttt{sigma88}), with \(\sigma^2 = I_8\).
\item
  \(I_n\): The \(n \times n\) identity matrix (\(n = 4, 8, 16\); the
  author's \texttt{ID4}, \texttt{ID8}, \texttt{ID16}).
\item
  \(0\): The zero matrix of the size required by its position (in a
  block matrix: a zero block).
\end{itemize}

(\texttt{B1\_stored\_eta4488} (passed).) The matrix \(\sigma\) has the
eigenvalues \(+1\) and \(-1\), four times each, so it is also a metric
of signature \((4,4)\).

\hypertarget{the-self-dual-and-anti-self-dual-4-times-4-matrices}{%
\subsection{\texorpdfstring{The self-dual and anti-self-dual
\(4 \times 4\)
matrices}{The self-dual and anti-self-dual 4 \textbackslash times 4 matrices}}\label{the-self-dual-and-anti-self-dual-4-times-4-matrices}}

The author's cell 24 defines, for \(h = 1, 2, 3\) and
\(p, q = 1, \dots, 4\),

\[
(s_h)_{pq} = \varepsilon_{hpq4} - \big(\delta_{p4}\delta_{qh} - \delta_{ph}\delta_{q4}\big),\qquad (t_h)_{pq} = \varepsilon_{hpq4} + \big(\delta_{p4}\delta_{qh} - \delta_{ph}\delta_{q4}\big)
\]

\textbf{Symbols in this equation.}

\begin{itemize}
\tightlist
\item
  \(s_h,\ t_h\): The author's \(4 \times 4\) matrices
  \texttt{s4by4{[}h{]}} and \texttt{t4by4{[}h{]}} (\(h = 1, 2, 3\)):
  real, antisymmetric, self-dual (\(s_h\)) and anti-self-dual (\(t_h\))
  in the four-dimensional sense.
\item
  \(h, p, q\): Integer labels of the notebook: \(h = 1, 2, 3\) numbers
  the matrices \(s_h\), \(t_h\); \(p, q = 1, \dots, 4\) are their row
  and column (counted from 1 as in Mathematica).
\item
  \(\varepsilon_{hpqr}\): The Levi-Civita symbol of four indices
  \(1, \dots, 4\) (the author's \texttt{Signature{[}\{h,p,q,r\}{]}}):
  \(+1\) for an even permutation of \((1, 2, 3, 4)\), \(-1\) for an odd
  one, \(0\) if an index repeats.
\item
  \(\delta_{pq}\): The Kronecker delta: \(1\) if \(p = q\), else \(0\)
  (the author's \texttt{ID4{[}{[}p,q{]}{]}}).
\end{itemize}

which gives (\texttt{B2\_stored\_s4by4\_t4by4} (passed))

\[
s_1 = \begin{pmatrix} 0 & 0 & 0 & 1 \\ 0 & 0 & 1 & 0 \\ 0 & -1 & 0 & 0 \\ -1 & 0 & 0 & 0 \end{pmatrix},\quad s_2 = \begin{pmatrix} 0 & 0 & -1 & 0 \\ 0 & 0 & 0 & 1 \\ 1 & 0 & 0 & 0 \\ 0 & -1 & 0 & 0 \end{pmatrix},\quad s_3 = \begin{pmatrix} 0 & 1 & 0 & 0 \\ -1 & 0 & 0 & 0 \\ 0 & 0 & 0 & 1 \\ 0 & 0 & -1 & 0 \end{pmatrix}
\]

\[
t_1 = \begin{pmatrix} 0 & 0 & 0 & -1 \\ 0 & 0 & 1 & 0 \\ 0 & -1 & 0 & 0 \\ 1 & 0 & 0 & 0 \end{pmatrix},\quad t_2 = \begin{pmatrix} 0 & 0 & -1 & 0 \\ 0 & 0 & 0 & -1 \\ 1 & 0 & 0 & 0 \\ 0 & 1 & 0 & 0 \end{pmatrix},\quad t_3 = \begin{pmatrix} 0 & 1 & 0 & 0 \\ -1 & 0 & 0 & 0 \\ 0 & 0 & 0 & -1 \\ 0 & 0 & 1 & 0 \end{pmatrix}
\]

\textbf{Symbols in this equation.}

\begin{itemize}
\tightlist
\item
  \(s_h,\ t_h\): The author's \(4 \times 4\) matrices
  \texttt{s4by4{[}h{]}} and \texttt{t4by4{[}h{]}} (\(h = 1, 2, 3\)):
  real, antisymmetric, self-dual (\(s_h\)) and anti-self-dual (\(t_h\))
  in the four-dimensional sense.
\end{itemize}

Each \(s_h\) equals its four-dimensional dual,
\(\tfrac12\sum_{j_1, j_2}\varepsilon_{pqj_1j_2}(s_h)_{j_1j_2} = (s_h)_{pq}\),
and each \(t_h\) minus its dual (the author's cells 67 and 70, both
giving zero matrices).

\hypertarget{the-8-times-8-matrices-tau_a-and-bartau_a}{%
\subsection{\texorpdfstring{The \(8 \times 8\) matrices \(\tau_a\) and
\(\bar\tau_a\)}{The 8 \textbackslash times 8 matrices \textbackslash tau\_a and \textbackslash bar\textbackslash tau\_a}}\label{the-8-times-8-matrices-tau_a-and-bartau_a}}

The author's cells 73 and 87 define

\[
\begin{aligned}
&\tau_0 = I_8,\qquad \tau_h = \begin{pmatrix} 0 & s_h \\ s_h & 0 \end{pmatrix},\qquad \tau_{7-h} = \begin{pmatrix} 0 & t_h \\ -t_h & 0 \end{pmatrix}\qquad (h = 1, 2, 3),\\
&\tau_7 = \tau_1\tau_2\tau_3\tau_4\tau_5\tau_6
\end{aligned}
\]

\textbf{Symbols in this equation.}

\begin{itemize}
\tightlist
\item
  \(\tau_a\): The author's real \(8 \times 8\) matrices
  \texttt{tau{[}a{]}} (\(a = 0, \dots, 7\)), \(\tau_a = \zeta_a\).
\item
  \(\zeta_A\): The author's intermediate \(8 \times 8\) matrices
  \texttt{zeta{[}A{]}} (\(A = 0, \dots, 7\)), built from \(s_h\) and
  \(t_h\).
\item
  \(s_h,\ t_h\): The author's \(4 \times 4\) matrices
  \texttt{s4by4{[}h{]}} and \texttt{t4by4{[}h{]}} (\(h = 1, 2, 3\)):
  real, antisymmetric, self-dual (\(s_h\)) and anti-self-dual (\(t_h\))
  in the four-dimensional sense.
\item
  \(I_n\): The \(n \times n\) identity matrix (\(n = 4, 8, 16\); the
  author's \texttt{ID4}, \texttt{ID8}, \texttt{ID16}).
\item
  \(0\): The zero matrix of the size required by its position (in a
  block matrix: a zero block).
\end{itemize}

(the author first names them \(\zeta_A\) and then sets
\(\tau_A = \zeta_A\)), and cell 91 defines

\[
\bar\tau_0 = I_8,\qquad \bar\tau_a = \sigma\,\tau_a^{\mathsf T}\,\sigma\quad (a = 1, \dots, 7)
\]

\textbf{Symbols in this equation.}

\begin{itemize}
\tightlist
\item
  \(\bar\tau_a\): The author's \texttt{taubar{[}a{]}} (written
  \(\overline{\tau}\) in the notebook), \(\bar\tau_0 = I_8\),
  \(\bar\tau_a = \sigma\,\tau_a^{\mathsf T}\,\sigma\) for
  \(a = 1, \dots, 7\) (not a complex conjugate).
\item
  \(\tau_a\): The author's real \(8 \times 8\) matrices
  \texttt{tau{[}a{]}} (\(a = 0, \dots, 7\)), \(\tau_a = \zeta_a\).
\item
  \(\sigma\): The real symmetric \(8 \times 8\) matrix
  \(\sigma = \begin{pmatrix} 0 & I_4 \\ I_4 & 0 \end{pmatrix}\) (the
  author's \texttt{sigma88}), with \(\sigma^2 = I_8\).
\item
  \(X^{\mathsf T},\ X^\dagger\): The transpose of a matrix and its
  Hermitian conjugate (transpose of the complex conjugate); for the real
  matrices of this document they coincide.
\item
  \(I_n\): The \(n \times n\) identity matrix (\(n = 4, 8, 16\); the
  author's \texttt{ID4}, \texttt{ID8}, \texttt{ID16}).
\end{itemize}

(\texttt{B4\_stored\_tau7} (passed), \texttt{C6\_taubar\_definition}
(passed)). Explicitly:

\[
\tau_0 = \begin{pmatrix} 1 & 0 & 0 & 0 & 0 & 0 & 0 & 0 \\ 0 & 1 & 0 & 0 & 0 & 0 & 0 & 0 \\ 0 & 0 & 1 & 0 & 0 & 0 & 0 & 0 \\ 0 & 0 & 0 & 1 & 0 & 0 & 0 & 0 \\ 0 & 0 & 0 & 0 & 1 & 0 & 0 & 0 \\ 0 & 0 & 0 & 0 & 0 & 1 & 0 & 0 \\ 0 & 0 & 0 & 0 & 0 & 0 & 1 & 0 \\ 0 & 0 & 0 & 0 & 0 & 0 & 0 & 1 \end{pmatrix},\qquad \bar\tau_0 = \begin{pmatrix} 1 & 0 & 0 & 0 & 0 & 0 & 0 & 0 \\ 0 & 1 & 0 & 0 & 0 & 0 & 0 & 0 \\ 0 & 0 & 1 & 0 & 0 & 0 & 0 & 0 \\ 0 & 0 & 0 & 1 & 0 & 0 & 0 & 0 \\ 0 & 0 & 0 & 0 & 1 & 0 & 0 & 0 \\ 0 & 0 & 0 & 0 & 0 & 1 & 0 & 0 \\ 0 & 0 & 0 & 0 & 0 & 0 & 1 & 0 \\ 0 & 0 & 0 & 0 & 0 & 0 & 0 & 1 \end{pmatrix}
\]

\[
\tau_1 = \begin{pmatrix} 0 & 0 & 0 & 0 & 0 & 0 & 0 & 1 \\ 0 & 0 & 0 & 0 & 0 & 0 & 1 & 0 \\ 0 & 0 & 0 & 0 & 0 & -1 & 0 & 0 \\ 0 & 0 & 0 & 0 & -1 & 0 & 0 & 0 \\ 0 & 0 & 0 & 1 & 0 & 0 & 0 & 0 \\ 0 & 0 & 1 & 0 & 0 & 0 & 0 & 0 \\ 0 & -1 & 0 & 0 & 0 & 0 & 0 & 0 \\ -1 & 0 & 0 & 0 & 0 & 0 & 0 & 0 \end{pmatrix},\qquad \bar\tau_1 = \begin{pmatrix} 0 & 0 & 0 & 0 & 0 & 0 & 0 & -1 \\ 0 & 0 & 0 & 0 & 0 & 0 & -1 & 0 \\ 0 & 0 & 0 & 0 & 0 & 1 & 0 & 0 \\ 0 & 0 & 0 & 0 & 1 & 0 & 0 & 0 \\ 0 & 0 & 0 & -1 & 0 & 0 & 0 & 0 \\ 0 & 0 & -1 & 0 & 0 & 0 & 0 & 0 \\ 0 & 1 & 0 & 0 & 0 & 0 & 0 & 0 \\ 1 & 0 & 0 & 0 & 0 & 0 & 0 & 0 \end{pmatrix}
\]

\[
\tau_2 = \begin{pmatrix} 0 & 0 & 0 & 0 & 0 & 0 & -1 & 0 \\ 0 & 0 & 0 & 0 & 0 & 0 & 0 & 1 \\ 0 & 0 & 0 & 0 & 1 & 0 & 0 & 0 \\ 0 & 0 & 0 & 0 & 0 & -1 & 0 & 0 \\ 0 & 0 & -1 & 0 & 0 & 0 & 0 & 0 \\ 0 & 0 & 0 & 1 & 0 & 0 & 0 & 0 \\ 1 & 0 & 0 & 0 & 0 & 0 & 0 & 0 \\ 0 & -1 & 0 & 0 & 0 & 0 & 0 & 0 \end{pmatrix},\qquad \bar\tau_2 = \begin{pmatrix} 0 & 0 & 0 & 0 & 0 & 0 & 1 & 0 \\ 0 & 0 & 0 & 0 & 0 & 0 & 0 & -1 \\ 0 & 0 & 0 & 0 & -1 & 0 & 0 & 0 \\ 0 & 0 & 0 & 0 & 0 & 1 & 0 & 0 \\ 0 & 0 & 1 & 0 & 0 & 0 & 0 & 0 \\ 0 & 0 & 0 & -1 & 0 & 0 & 0 & 0 \\ -1 & 0 & 0 & 0 & 0 & 0 & 0 & 0 \\ 0 & 1 & 0 & 0 & 0 & 0 & 0 & 0 \end{pmatrix}
\]

\[
\tau_3 = \begin{pmatrix} 0 & 0 & 0 & 0 & 0 & 1 & 0 & 0 \\ 0 & 0 & 0 & 0 & -1 & 0 & 0 & 0 \\ 0 & 0 & 0 & 0 & 0 & 0 & 0 & 1 \\ 0 & 0 & 0 & 0 & 0 & 0 & -1 & 0 \\ 0 & 1 & 0 & 0 & 0 & 0 & 0 & 0 \\ -1 & 0 & 0 & 0 & 0 & 0 & 0 & 0 \\ 0 & 0 & 0 & 1 & 0 & 0 & 0 & 0 \\ 0 & 0 & -1 & 0 & 0 & 0 & 0 & 0 \end{pmatrix},\qquad \bar\tau_3 = \begin{pmatrix} 0 & 0 & 0 & 0 & 0 & -1 & 0 & 0 \\ 0 & 0 & 0 & 0 & 1 & 0 & 0 & 0 \\ 0 & 0 & 0 & 0 & 0 & 0 & 0 & -1 \\ 0 & 0 & 0 & 0 & 0 & 0 & 1 & 0 \\ 0 & -1 & 0 & 0 & 0 & 0 & 0 & 0 \\ 1 & 0 & 0 & 0 & 0 & 0 & 0 & 0 \\ 0 & 0 & 0 & -1 & 0 & 0 & 0 & 0 \\ 0 & 0 & 1 & 0 & 0 & 0 & 0 & 0 \end{pmatrix}
\]

\[
\tau_4 = \begin{pmatrix} 0 & 0 & 0 & 0 & 0 & 1 & 0 & 0 \\ 0 & 0 & 0 & 0 & -1 & 0 & 0 & 0 \\ 0 & 0 & 0 & 0 & 0 & 0 & 0 & -1 \\ 0 & 0 & 0 & 0 & 0 & 0 & 1 & 0 \\ 0 & -1 & 0 & 0 & 0 & 0 & 0 & 0 \\ 1 & 0 & 0 & 0 & 0 & 0 & 0 & 0 \\ 0 & 0 & 0 & 1 & 0 & 0 & 0 & 0 \\ 0 & 0 & -1 & 0 & 0 & 0 & 0 & 0 \end{pmatrix},\qquad \bar\tau_4 = \begin{pmatrix} 0 & 0 & 0 & 0 & 0 & -1 & 0 & 0 \\ 0 & 0 & 0 & 0 & 1 & 0 & 0 & 0 \\ 0 & 0 & 0 & 0 & 0 & 0 & 0 & 1 \\ 0 & 0 & 0 & 0 & 0 & 0 & -1 & 0 \\ 0 & 1 & 0 & 0 & 0 & 0 & 0 & 0 \\ -1 & 0 & 0 & 0 & 0 & 0 & 0 & 0 \\ 0 & 0 & 0 & -1 & 0 & 0 & 0 & 0 \\ 0 & 0 & 1 & 0 & 0 & 0 & 0 & 0 \end{pmatrix}
\]

\[
\tau_5 = \begin{pmatrix} 0 & 0 & 0 & 0 & 0 & 0 & -1 & 0 \\ 0 & 0 & 0 & 0 & 0 & 0 & 0 & -1 \\ 0 & 0 & 0 & 0 & 1 & 0 & 0 & 0 \\ 0 & 0 & 0 & 0 & 0 & 1 & 0 & 0 \\ 0 & 0 & 1 & 0 & 0 & 0 & 0 & 0 \\ 0 & 0 & 0 & 1 & 0 & 0 & 0 & 0 \\ -1 & 0 & 0 & 0 & 0 & 0 & 0 & 0 \\ 0 & -1 & 0 & 0 & 0 & 0 & 0 & 0 \end{pmatrix},\qquad \bar\tau_5 = \begin{pmatrix} 0 & 0 & 0 & 0 & 0 & 0 & 1 & 0 \\ 0 & 0 & 0 & 0 & 0 & 0 & 0 & 1 \\ 0 & 0 & 0 & 0 & -1 & 0 & 0 & 0 \\ 0 & 0 & 0 & 0 & 0 & -1 & 0 & 0 \\ 0 & 0 & -1 & 0 & 0 & 0 & 0 & 0 \\ 0 & 0 & 0 & -1 & 0 & 0 & 0 & 0 \\ 1 & 0 & 0 & 0 & 0 & 0 & 0 & 0 \\ 0 & 1 & 0 & 0 & 0 & 0 & 0 & 0 \end{pmatrix}
\]

\[
\tau_6 = \begin{pmatrix} 0 & 0 & 0 & 0 & 0 & 0 & 0 & -1 \\ 0 & 0 & 0 & 0 & 0 & 0 & 1 & 0 \\ 0 & 0 & 0 & 0 & 0 & -1 & 0 & 0 \\ 0 & 0 & 0 & 0 & 1 & 0 & 0 & 0 \\ 0 & 0 & 0 & 1 & 0 & 0 & 0 & 0 \\ 0 & 0 & -1 & 0 & 0 & 0 & 0 & 0 \\ 0 & 1 & 0 & 0 & 0 & 0 & 0 & 0 \\ -1 & 0 & 0 & 0 & 0 & 0 & 0 & 0 \end{pmatrix},\qquad \bar\tau_6 = \begin{pmatrix} 0 & 0 & 0 & 0 & 0 & 0 & 0 & 1 \\ 0 & 0 & 0 & 0 & 0 & 0 & -1 & 0 \\ 0 & 0 & 0 & 0 & 0 & 1 & 0 & 0 \\ 0 & 0 & 0 & 0 & -1 & 0 & 0 & 0 \\ 0 & 0 & 0 & -1 & 0 & 0 & 0 & 0 \\ 0 & 0 & 1 & 0 & 0 & 0 & 0 & 0 \\ 0 & -1 & 0 & 0 & 0 & 0 & 0 & 0 \\ 1 & 0 & 0 & 0 & 0 & 0 & 0 & 0 \end{pmatrix}
\]

\[
\tau_7 = \begin{pmatrix} -1 & 0 & 0 & 0 & 0 & 0 & 0 & 0 \\ 0 & -1 & 0 & 0 & 0 & 0 & 0 & 0 \\ 0 & 0 & -1 & 0 & 0 & 0 & 0 & 0 \\ 0 & 0 & 0 & -1 & 0 & 0 & 0 & 0 \\ 0 & 0 & 0 & 0 & 1 & 0 & 0 & 0 \\ 0 & 0 & 0 & 0 & 0 & 1 & 0 & 0 \\ 0 & 0 & 0 & 0 & 0 & 0 & 1 & 0 \\ 0 & 0 & 0 & 0 & 0 & 0 & 0 & 1 \end{pmatrix},\qquad \bar\tau_7 = \begin{pmatrix} 1 & 0 & 0 & 0 & 0 & 0 & 0 & 0 \\ 0 & 1 & 0 & 0 & 0 & 0 & 0 & 0 \\ 0 & 0 & 1 & 0 & 0 & 0 & 0 & 0 \\ 0 & 0 & 0 & 1 & 0 & 0 & 0 & 0 \\ 0 & 0 & 0 & 0 & -1 & 0 & 0 & 0 \\ 0 & 0 & 0 & 0 & 0 & -1 & 0 & 0 \\ 0 & 0 & 0 & 0 & 0 & 0 & -1 & 0 \\ 0 & 0 & 0 & 0 & 0 & 0 & 0 & -1 \end{pmatrix}
\]

\textbf{Symbols in this equation.}

\begin{itemize}
\tightlist
\item
  \(\tau_a\): The author's real \(8 \times 8\) matrices
  \texttt{tau{[}a{]}} (\(a = 0, \dots, 7\)), \(\tau_a = \zeta_a\).
\item
  \(\bar\tau_a\): The author's \texttt{taubar{[}a{]}} (written
  \(\overline{\tau}\) in the notebook), \(\bar\tau_0 = I_8\),
  \(\bar\tau_a = \sigma\,\tau_a^{\mathsf T}\,\sigma\) for
  \(a = 1, \dots, 7\) (not a complex conjugate).
\end{itemize}

They satisfy, for all \(a, b = 0, \dots, 7\)
(\texttt{C5\_block\_relations} (passed)),

\[
\tau_a\bar\tau_b + \tau_b\bar\tau_a = 2\eta_{ab}I_8,\qquad \bar\tau_a\tau_b + \bar\tau_b\tau_a = 2\eta_{ab}I_8
\]

\textbf{Symbols in this equation.}

\begin{itemize}
\tightlist
\item
  \(\tau_a\): The author's real \(8 \times 8\) matrices
  \texttt{tau{[}a{]}} (\(a = 0, \dots, 7\)), \(\tau_a = \zeta_a\).
\item
  \(\bar\tau_a\): The author's \texttt{taubar{[}a{]}} (written
  \(\overline{\tau}\) in the notebook), \(\bar\tau_0 = I_8\),
  \(\bar\tau_a = \sigma\,\tau_a^{\mathsf T}\,\sigma\) for
  \(a = 1, \dots, 7\) (not a complex conjugate).
\item
  \(\eta_{ab}\): The flat frame metric,
  \(\eta = \mathrm{diag}(+1, +1, +1, +1, -1, -1, -1, -1)\) in the order
  \(x_0, \dots, x_7\) (the author's \texttt{eta4488}); \(\eta^{ab}\) is
  the same matrix.
\item
  \(I_n\): The \(n \times n\) identity matrix (\(n = 4, 8, 16\); the
  author's \texttt{ID4}, \texttt{ID8}, \texttt{ID16}).
\end{itemize}

and each \(\tau_a\) is orthogonal, \(\tau_a\tau_a^{\mathsf T} = I_8\)
(\texttt{C7\_tau\_orthogonal} (passed)).

\hypertarget{the-eight-16-times-16-dirac-matrices}{%
\subsection{\texorpdfstring{The eight \(16 \times 16\) Dirac
matrices}{The eight 16 \textbackslash times 16 Dirac matrices}}\label{the-eight-16-times-16-dirac-matrices}}

The author's cell 133 (\texttt{In{[}97{]}}) defines, for
\(a = 0, \dots, 7\),

\[
T_a = \begin{pmatrix} 0 & \bar\tau_a \\ \tau_a & 0 \end{pmatrix}
\]

\textbf{Symbols in this equation.}

\begin{itemize}
\tightlist
\item
  \(T_a\): The author's eight real \(16 \times 16\) Dirac matrices
  \texttt{T16\^{}A{[}a{]}} (\(a = 0, \dots, 7\)), the gamma matrix of
  the frame direction of \(x_a\): \(\gamma^{(x_a)} = T_a\).
\item
  \(\tau_a\): The author's real \(8 \times 8\) matrices
  \texttt{tau{[}a{]}} (\(a = 0, \dots, 7\)), \(\tau_a = \zeta_a\).
\item
  \(\bar\tau_a\): The author's \texttt{taubar{[}a{]}} (written
  \(\overline{\tau}\) in the notebook), \(\bar\tau_0 = I_8\),
  \(\bar\tau_a = \sigma\,\tau_a^{\mathsf T}\,\sigma\) for
  \(a = 1, \dots, 7\) (not a complex conjugate).
\item
  \(0\): The zero matrix of the size required by its position (in a
  block matrix: a zero block).
\end{itemize}

Here they are, every entry (\texttt{B11\_stored\_T16A\_0\_to\_7}
(passed)). Below each matrix, its action as a signed permutation is
listed in the form ``row:sign column'' (the only non-zero entry of each
row; for example 0:+8 means that row 0 has the entry \(+1\) in column
8).

\textbf{\(T_0\) = \texttt{T16\^{}A{[}0{]}}, the Dirac matrix of the
direction \(x_0\) (the hidden space-like direction);
\(T_0^2 = I_{16}\).}

\[
T_0 = \scriptstyle \begin{pmatrix} 0 & 0 & 0 & 0 & 0 & 0 & 0 & 0 & 1 & 0 & 0 & 0 & 0 & 0 & 0 & 0 \\ 0 & 0 & 0 & 0 & 0 & 0 & 0 & 0 & 0 & 1 & 0 & 0 & 0 & 0 & 0 & 0 \\ 0 & 0 & 0 & 0 & 0 & 0 & 0 & 0 & 0 & 0 & 1 & 0 & 0 & 0 & 0 & 0 \\ 0 & 0 & 0 & 0 & 0 & 0 & 0 & 0 & 0 & 0 & 0 & 1 & 0 & 0 & 0 & 0 \\ 0 & 0 & 0 & 0 & 0 & 0 & 0 & 0 & 0 & 0 & 0 & 0 & 1 & 0 & 0 & 0 \\ 0 & 0 & 0 & 0 & 0 & 0 & 0 & 0 & 0 & 0 & 0 & 0 & 0 & 1 & 0 & 0 \\ 0 & 0 & 0 & 0 & 0 & 0 & 0 & 0 & 0 & 0 & 0 & 0 & 0 & 0 & 1 & 0 \\ 0 & 0 & 0 & 0 & 0 & 0 & 0 & 0 & 0 & 0 & 0 & 0 & 0 & 0 & 0 & 1 \\ 1 & 0 & 0 & 0 & 0 & 0 & 0 & 0 & 0 & 0 & 0 & 0 & 0 & 0 & 0 & 0 \\ 0 & 1 & 0 & 0 & 0 & 0 & 0 & 0 & 0 & 0 & 0 & 0 & 0 & 0 & 0 & 0 \\ 0 & 0 & 1 & 0 & 0 & 0 & 0 & 0 & 0 & 0 & 0 & 0 & 0 & 0 & 0 & 0 \\ 0 & 0 & 0 & 1 & 0 & 0 & 0 & 0 & 0 & 0 & 0 & 0 & 0 & 0 & 0 & 0 \\ 0 & 0 & 0 & 0 & 1 & 0 & 0 & 0 & 0 & 0 & 0 & 0 & 0 & 0 & 0 & 0 \\ 0 & 0 & 0 & 0 & 0 & 1 & 0 & 0 & 0 & 0 & 0 & 0 & 0 & 0 & 0 & 0 \\ 0 & 0 & 0 & 0 & 0 & 0 & 1 & 0 & 0 & 0 & 0 & 0 & 0 & 0 & 0 & 0 \\ 0 & 0 & 0 & 0 & 0 & 0 & 0 & 1 & 0 & 0 & 0 & 0 & 0 & 0 & 0 & 0 \end{pmatrix}
\]

Signed permutation of \(T_0\): 0:+8, 1:+9, 2:+10, 3:+11, 4:+12, 5:+13,
6:+14, 7:+15, 8:+0, 9:+1, 10:+2, 11:+3, 12:+4, 13:+5, 14:+6, 15:+7.

\textbf{\(T_1\) = \texttt{T16\^{}A{[}1{]}}, the Dirac matrix of the
direction \(x_1\) (ordinary space); \(T_1^2 = I_{16}\).}

\[
T_1 = \scriptstyle \begin{pmatrix} 0 & 0 & 0 & 0 & 0 & 0 & 0 & 0 & 0 & 0 & 0 & 0 & 0 & 0 & 0 & -1 \\ 0 & 0 & 0 & 0 & 0 & 0 & 0 & 0 & 0 & 0 & 0 & 0 & 0 & 0 & -1 & 0 \\ 0 & 0 & 0 & 0 & 0 & 0 & 0 & 0 & 0 & 0 & 0 & 0 & 0 & 1 & 0 & 0 \\ 0 & 0 & 0 & 0 & 0 & 0 & 0 & 0 & 0 & 0 & 0 & 0 & 1 & 0 & 0 & 0 \\ 0 & 0 & 0 & 0 & 0 & 0 & 0 & 0 & 0 & 0 & 0 & -1 & 0 & 0 & 0 & 0 \\ 0 & 0 & 0 & 0 & 0 & 0 & 0 & 0 & 0 & 0 & -1 & 0 & 0 & 0 & 0 & 0 \\ 0 & 0 & 0 & 0 & 0 & 0 & 0 & 0 & 0 & 1 & 0 & 0 & 0 & 0 & 0 & 0 \\ 0 & 0 & 0 & 0 & 0 & 0 & 0 & 0 & 1 & 0 & 0 & 0 & 0 & 0 & 0 & 0 \\ 0 & 0 & 0 & 0 & 0 & 0 & 0 & 1 & 0 & 0 & 0 & 0 & 0 & 0 & 0 & 0 \\ 0 & 0 & 0 & 0 & 0 & 0 & 1 & 0 & 0 & 0 & 0 & 0 & 0 & 0 & 0 & 0 \\ 0 & 0 & 0 & 0 & 0 & -1 & 0 & 0 & 0 & 0 & 0 & 0 & 0 & 0 & 0 & 0 \\ 0 & 0 & 0 & 0 & -1 & 0 & 0 & 0 & 0 & 0 & 0 & 0 & 0 & 0 & 0 & 0 \\ 0 & 0 & 0 & 1 & 0 & 0 & 0 & 0 & 0 & 0 & 0 & 0 & 0 & 0 & 0 & 0 \\ 0 & 0 & 1 & 0 & 0 & 0 & 0 & 0 & 0 & 0 & 0 & 0 & 0 & 0 & 0 & 0 \\ 0 & -1 & 0 & 0 & 0 & 0 & 0 & 0 & 0 & 0 & 0 & 0 & 0 & 0 & 0 & 0 \\ -1 & 0 & 0 & 0 & 0 & 0 & 0 & 0 & 0 & 0 & 0 & 0 & 0 & 0 & 0 & 0 \end{pmatrix}
\]

Signed permutation of \(T_1\): 0:-15, 1:-14, 2:+13, 3:+12, 4:-11, 5:-10,
6:+9, 7:+8, 8:+7, 9:+6, 10:-5, 11:-4, 12:+3, 13:+2, 14:-1, 15:-0.

\textbf{\(T_2\) = \texttt{T16\^{}A{[}2{]}}, the Dirac matrix of the
direction \(x_2\) (ordinary space); \(T_2^2 = I_{16}\).}

\[
T_2 = \scriptstyle \begin{pmatrix} 0 & 0 & 0 & 0 & 0 & 0 & 0 & 0 & 0 & 0 & 0 & 0 & 0 & 0 & 1 & 0 \\ 0 & 0 & 0 & 0 & 0 & 0 & 0 & 0 & 0 & 0 & 0 & 0 & 0 & 0 & 0 & -1 \\ 0 & 0 & 0 & 0 & 0 & 0 & 0 & 0 & 0 & 0 & 0 & 0 & -1 & 0 & 0 & 0 \\ 0 & 0 & 0 & 0 & 0 & 0 & 0 & 0 & 0 & 0 & 0 & 0 & 0 & 1 & 0 & 0 \\ 0 & 0 & 0 & 0 & 0 & 0 & 0 & 0 & 0 & 0 & 1 & 0 & 0 & 0 & 0 & 0 \\ 0 & 0 & 0 & 0 & 0 & 0 & 0 & 0 & 0 & 0 & 0 & -1 & 0 & 0 & 0 & 0 \\ 0 & 0 & 0 & 0 & 0 & 0 & 0 & 0 & -1 & 0 & 0 & 0 & 0 & 0 & 0 & 0 \\ 0 & 0 & 0 & 0 & 0 & 0 & 0 & 0 & 0 & 1 & 0 & 0 & 0 & 0 & 0 & 0 \\ 0 & 0 & 0 & 0 & 0 & 0 & -1 & 0 & 0 & 0 & 0 & 0 & 0 & 0 & 0 & 0 \\ 0 & 0 & 0 & 0 & 0 & 0 & 0 & 1 & 0 & 0 & 0 & 0 & 0 & 0 & 0 & 0 \\ 0 & 0 & 0 & 0 & 1 & 0 & 0 & 0 & 0 & 0 & 0 & 0 & 0 & 0 & 0 & 0 \\ 0 & 0 & 0 & 0 & 0 & -1 & 0 & 0 & 0 & 0 & 0 & 0 & 0 & 0 & 0 & 0 \\ 0 & 0 & -1 & 0 & 0 & 0 & 0 & 0 & 0 & 0 & 0 & 0 & 0 & 0 & 0 & 0 \\ 0 & 0 & 0 & 1 & 0 & 0 & 0 & 0 & 0 & 0 & 0 & 0 & 0 & 0 & 0 & 0 \\ 1 & 0 & 0 & 0 & 0 & 0 & 0 & 0 & 0 & 0 & 0 & 0 & 0 & 0 & 0 & 0 \\ 0 & -1 & 0 & 0 & 0 & 0 & 0 & 0 & 0 & 0 & 0 & 0 & 0 & 0 & 0 & 0 \end{pmatrix}
\]

Signed permutation of \(T_2\): 0:+14, 1:-15, 2:-12, 3:+13, 4:+10, 5:-11,
6:-8, 7:+9, 8:-6, 9:+7, 10:+4, 11:-5, 12:-2, 13:+3, 14:+0, 15:-1.

\textbf{\(T_3\) = \texttt{T16\^{}A{[}3{]}}, the Dirac matrix of the
direction \(x_3\) (ordinary space); \(T_3^2 = I_{16}\).}

\[
T_3 = \scriptstyle \begin{pmatrix} 0 & 0 & 0 & 0 & 0 & 0 & 0 & 0 & 0 & 0 & 0 & 0 & 0 & -1 & 0 & 0 \\ 0 & 0 & 0 & 0 & 0 & 0 & 0 & 0 & 0 & 0 & 0 & 0 & 1 & 0 & 0 & 0 \\ 0 & 0 & 0 & 0 & 0 & 0 & 0 & 0 & 0 & 0 & 0 & 0 & 0 & 0 & 0 & -1 \\ 0 & 0 & 0 & 0 & 0 & 0 & 0 & 0 & 0 & 0 & 0 & 0 & 0 & 0 & 1 & 0 \\ 0 & 0 & 0 & 0 & 0 & 0 & 0 & 0 & 0 & -1 & 0 & 0 & 0 & 0 & 0 & 0 \\ 0 & 0 & 0 & 0 & 0 & 0 & 0 & 0 & 1 & 0 & 0 & 0 & 0 & 0 & 0 & 0 \\ 0 & 0 & 0 & 0 & 0 & 0 & 0 & 0 & 0 & 0 & 0 & -1 & 0 & 0 & 0 & 0 \\ 0 & 0 & 0 & 0 & 0 & 0 & 0 & 0 & 0 & 0 & 1 & 0 & 0 & 0 & 0 & 0 \\ 0 & 0 & 0 & 0 & 0 & 1 & 0 & 0 & 0 & 0 & 0 & 0 & 0 & 0 & 0 & 0 \\ 0 & 0 & 0 & 0 & -1 & 0 & 0 & 0 & 0 & 0 & 0 & 0 & 0 & 0 & 0 & 0 \\ 0 & 0 & 0 & 0 & 0 & 0 & 0 & 1 & 0 & 0 & 0 & 0 & 0 & 0 & 0 & 0 \\ 0 & 0 & 0 & 0 & 0 & 0 & -1 & 0 & 0 & 0 & 0 & 0 & 0 & 0 & 0 & 0 \\ 0 & 1 & 0 & 0 & 0 & 0 & 0 & 0 & 0 & 0 & 0 & 0 & 0 & 0 & 0 & 0 \\ -1 & 0 & 0 & 0 & 0 & 0 & 0 & 0 & 0 & 0 & 0 & 0 & 0 & 0 & 0 & 0 \\ 0 & 0 & 0 & 1 & 0 & 0 & 0 & 0 & 0 & 0 & 0 & 0 & 0 & 0 & 0 & 0 \\ 0 & 0 & -1 & 0 & 0 & 0 & 0 & 0 & 0 & 0 & 0 & 0 & 0 & 0 & 0 & 0 \end{pmatrix}
\]

Signed permutation of \(T_3\): 0:-13, 1:+12, 2:-15, 3:+14, 4:-9, 5:+8,
6:-11, 7:+10, 8:+5, 9:-4, 10:+7, 11:-6, 12:+1, 13:-0, 14:+3, 15:-2.

\textbf{\(T_4\) = \texttt{T16\^{}A{[}4{]}}, the Dirac matrix of the
direction \(x_4\) (the time); \(T_4^2 = -I_{16}\).}

\[
T_4 = \scriptstyle \begin{pmatrix} 0 & 0 & 0 & 0 & 0 & 0 & 0 & 0 & 0 & 0 & 0 & 0 & 0 & -1 & 0 & 0 \\ 0 & 0 & 0 & 0 & 0 & 0 & 0 & 0 & 0 & 0 & 0 & 0 & 1 & 0 & 0 & 0 \\ 0 & 0 & 0 & 0 & 0 & 0 & 0 & 0 & 0 & 0 & 0 & 0 & 0 & 0 & 0 & 1 \\ 0 & 0 & 0 & 0 & 0 & 0 & 0 & 0 & 0 & 0 & 0 & 0 & 0 & 0 & -1 & 0 \\ 0 & 0 & 0 & 0 & 0 & 0 & 0 & 0 & 0 & 1 & 0 & 0 & 0 & 0 & 0 & 0 \\ 0 & 0 & 0 & 0 & 0 & 0 & 0 & 0 & -1 & 0 & 0 & 0 & 0 & 0 & 0 & 0 \\ 0 & 0 & 0 & 0 & 0 & 0 & 0 & 0 & 0 & 0 & 0 & -1 & 0 & 0 & 0 & 0 \\ 0 & 0 & 0 & 0 & 0 & 0 & 0 & 0 & 0 & 0 & 1 & 0 & 0 & 0 & 0 & 0 \\ 0 & 0 & 0 & 0 & 0 & 1 & 0 & 0 & 0 & 0 & 0 & 0 & 0 & 0 & 0 & 0 \\ 0 & 0 & 0 & 0 & -1 & 0 & 0 & 0 & 0 & 0 & 0 & 0 & 0 & 0 & 0 & 0 \\ 0 & 0 & 0 & 0 & 0 & 0 & 0 & -1 & 0 & 0 & 0 & 0 & 0 & 0 & 0 & 0 \\ 0 & 0 & 0 & 0 & 0 & 0 & 1 & 0 & 0 & 0 & 0 & 0 & 0 & 0 & 0 & 0 \\ 0 & -1 & 0 & 0 & 0 & 0 & 0 & 0 & 0 & 0 & 0 & 0 & 0 & 0 & 0 & 0 \\ 1 & 0 & 0 & 0 & 0 & 0 & 0 & 0 & 0 & 0 & 0 & 0 & 0 & 0 & 0 & 0 \\ 0 & 0 & 0 & 1 & 0 & 0 & 0 & 0 & 0 & 0 & 0 & 0 & 0 & 0 & 0 & 0 \\ 0 & 0 & -1 & 0 & 0 & 0 & 0 & 0 & 0 & 0 & 0 & 0 & 0 & 0 & 0 & 0 \end{pmatrix}
\]

Signed permutation of \(T_4\): 0:-13, 1:+12, 2:+15, 3:-14, 4:+9, 5:-8,
6:-11, 7:+10, 8:+5, 9:-4, 10:-7, 11:+6, 12:-1, 13:+0, 14:+3, 15:-2.

\textbf{\(T_5\) = \texttt{T16\^{}A{[}5{]}}, the Dirac matrix of the
direction \(x_5\) (an extra time-like direction); \(T_5^2 = -I_{16}\).}

\[
T_5 = \scriptstyle \begin{pmatrix} 0 & 0 & 0 & 0 & 0 & 0 & 0 & 0 & 0 & 0 & 0 & 0 & 0 & 0 & 1 & 0 \\ 0 & 0 & 0 & 0 & 0 & 0 & 0 & 0 & 0 & 0 & 0 & 0 & 0 & 0 & 0 & 1 \\ 0 & 0 & 0 & 0 & 0 & 0 & 0 & 0 & 0 & 0 & 0 & 0 & -1 & 0 & 0 & 0 \\ 0 & 0 & 0 & 0 & 0 & 0 & 0 & 0 & 0 & 0 & 0 & 0 & 0 & -1 & 0 & 0 \\ 0 & 0 & 0 & 0 & 0 & 0 & 0 & 0 & 0 & 0 & -1 & 0 & 0 & 0 & 0 & 0 \\ 0 & 0 & 0 & 0 & 0 & 0 & 0 & 0 & 0 & 0 & 0 & -1 & 0 & 0 & 0 & 0 \\ 0 & 0 & 0 & 0 & 0 & 0 & 0 & 0 & 1 & 0 & 0 & 0 & 0 & 0 & 0 & 0 \\ 0 & 0 & 0 & 0 & 0 & 0 & 0 & 0 & 0 & 1 & 0 & 0 & 0 & 0 & 0 & 0 \\ 0 & 0 & 0 & 0 & 0 & 0 & -1 & 0 & 0 & 0 & 0 & 0 & 0 & 0 & 0 & 0 \\ 0 & 0 & 0 & 0 & 0 & 0 & 0 & -1 & 0 & 0 & 0 & 0 & 0 & 0 & 0 & 0 \\ 0 & 0 & 0 & 0 & 1 & 0 & 0 & 0 & 0 & 0 & 0 & 0 & 0 & 0 & 0 & 0 \\ 0 & 0 & 0 & 0 & 0 & 1 & 0 & 0 & 0 & 0 & 0 & 0 & 0 & 0 & 0 & 0 \\ 0 & 0 & 1 & 0 & 0 & 0 & 0 & 0 & 0 & 0 & 0 & 0 & 0 & 0 & 0 & 0 \\ 0 & 0 & 0 & 1 & 0 & 0 & 0 & 0 & 0 & 0 & 0 & 0 & 0 & 0 & 0 & 0 \\ -1 & 0 & 0 & 0 & 0 & 0 & 0 & 0 & 0 & 0 & 0 & 0 & 0 & 0 & 0 & 0 \\ 0 & -1 & 0 & 0 & 0 & 0 & 0 & 0 & 0 & 0 & 0 & 0 & 0 & 0 & 0 & 0 \end{pmatrix}
\]

Signed permutation of \(T_5\): 0:+14, 1:+15, 2:-12, 3:-13, 4:-10, 5:-11,
6:+8, 7:+9, 8:-6, 9:-7, 10:+4, 11:+5, 12:+2, 13:+3, 14:-0, 15:-1.

\textbf{\(T_6\) = \texttt{T16\^{}A{[}6{]}}, the Dirac matrix of the
direction \(x_6\) (an extra time-like direction); \(T_6^2 = -I_{16}\).}

\[
T_6 = \scriptstyle \begin{pmatrix} 0 & 0 & 0 & 0 & 0 & 0 & 0 & 0 & 0 & 0 & 0 & 0 & 0 & 0 & 0 & 1 \\ 0 & 0 & 0 & 0 & 0 & 0 & 0 & 0 & 0 & 0 & 0 & 0 & 0 & 0 & -1 & 0 \\ 0 & 0 & 0 & 0 & 0 & 0 & 0 & 0 & 0 & 0 & 0 & 0 & 0 & 1 & 0 & 0 \\ 0 & 0 & 0 & 0 & 0 & 0 & 0 & 0 & 0 & 0 & 0 & 0 & -1 & 0 & 0 & 0 \\ 0 & 0 & 0 & 0 & 0 & 0 & 0 & 0 & 0 & 0 & 0 & -1 & 0 & 0 & 0 & 0 \\ 0 & 0 & 0 & 0 & 0 & 0 & 0 & 0 & 0 & 0 & 1 & 0 & 0 & 0 & 0 & 0 \\ 0 & 0 & 0 & 0 & 0 & 0 & 0 & 0 & 0 & -1 & 0 & 0 & 0 & 0 & 0 & 0 \\ 0 & 0 & 0 & 0 & 0 & 0 & 0 & 0 & 1 & 0 & 0 & 0 & 0 & 0 & 0 & 0 \\ 0 & 0 & 0 & 0 & 0 & 0 & 0 & -1 & 0 & 0 & 0 & 0 & 0 & 0 & 0 & 0 \\ 0 & 0 & 0 & 0 & 0 & 0 & 1 & 0 & 0 & 0 & 0 & 0 & 0 & 0 & 0 & 0 \\ 0 & 0 & 0 & 0 & 0 & -1 & 0 & 0 & 0 & 0 & 0 & 0 & 0 & 0 & 0 & 0 \\ 0 & 0 & 0 & 0 & 1 & 0 & 0 & 0 & 0 & 0 & 0 & 0 & 0 & 0 & 0 & 0 \\ 0 & 0 & 0 & 1 & 0 & 0 & 0 & 0 & 0 & 0 & 0 & 0 & 0 & 0 & 0 & 0 \\ 0 & 0 & -1 & 0 & 0 & 0 & 0 & 0 & 0 & 0 & 0 & 0 & 0 & 0 & 0 & 0 \\ 0 & 1 & 0 & 0 & 0 & 0 & 0 & 0 & 0 & 0 & 0 & 0 & 0 & 0 & 0 & 0 \\ -1 & 0 & 0 & 0 & 0 & 0 & 0 & 0 & 0 & 0 & 0 & 0 & 0 & 0 & 0 & 0 \end{pmatrix}
\]

Signed permutation of \(T_6\): 0:+15, 1:-14, 2:+13, 3:-12, 4:-11, 5:+10,
6:-9, 7:+8, 8:-7, 9:+6, 10:-5, 11:+4, 12:+3, 13:-2, 14:+1, 15:-0.

\textbf{\(T_7\) = \texttt{T16\^{}A{[}7{]}}, the Dirac matrix of the
direction \(x_7\) (an extra time-like direction); \(T_7^2 = -I_{16}\).}

\[
T_7 = \scriptstyle \begin{pmatrix} 0 & 0 & 0 & 0 & 0 & 0 & 0 & 0 & 1 & 0 & 0 & 0 & 0 & 0 & 0 & 0 \\ 0 & 0 & 0 & 0 & 0 & 0 & 0 & 0 & 0 & 1 & 0 & 0 & 0 & 0 & 0 & 0 \\ 0 & 0 & 0 & 0 & 0 & 0 & 0 & 0 & 0 & 0 & 1 & 0 & 0 & 0 & 0 & 0 \\ 0 & 0 & 0 & 0 & 0 & 0 & 0 & 0 & 0 & 0 & 0 & 1 & 0 & 0 & 0 & 0 \\ 0 & 0 & 0 & 0 & 0 & 0 & 0 & 0 & 0 & 0 & 0 & 0 & -1 & 0 & 0 & 0 \\ 0 & 0 & 0 & 0 & 0 & 0 & 0 & 0 & 0 & 0 & 0 & 0 & 0 & -1 & 0 & 0 \\ 0 & 0 & 0 & 0 & 0 & 0 & 0 & 0 & 0 & 0 & 0 & 0 & 0 & 0 & -1 & 0 \\ 0 & 0 & 0 & 0 & 0 & 0 & 0 & 0 & 0 & 0 & 0 & 0 & 0 & 0 & 0 & -1 \\ -1 & 0 & 0 & 0 & 0 & 0 & 0 & 0 & 0 & 0 & 0 & 0 & 0 & 0 & 0 & 0 \\ 0 & -1 & 0 & 0 & 0 & 0 & 0 & 0 & 0 & 0 & 0 & 0 & 0 & 0 & 0 & 0 \\ 0 & 0 & -1 & 0 & 0 & 0 & 0 & 0 & 0 & 0 & 0 & 0 & 0 & 0 & 0 & 0 \\ 0 & 0 & 0 & -1 & 0 & 0 & 0 & 0 & 0 & 0 & 0 & 0 & 0 & 0 & 0 & 0 \\ 0 & 0 & 0 & 0 & 1 & 0 & 0 & 0 & 0 & 0 & 0 & 0 & 0 & 0 & 0 & 0 \\ 0 & 0 & 0 & 0 & 0 & 1 & 0 & 0 & 0 & 0 & 0 & 0 & 0 & 0 & 0 & 0 \\ 0 & 0 & 0 & 0 & 0 & 0 & 1 & 0 & 0 & 0 & 0 & 0 & 0 & 0 & 0 & 0 \\ 0 & 0 & 0 & 0 & 0 & 0 & 0 & 1 & 0 & 0 & 0 & 0 & 0 & 0 & 0 & 0 \end{pmatrix}
\]

Signed permutation of \(T_7\): 0:+8, 1:+9, 2:+10, 3:+11, 4:-12, 5:-13,
6:-14, 7:-15, 8:-0, 9:-1, 10:-2, 11:-3, 12:+4, 13:+5, 14:+6, 15:+7.

\textbf{Symbols in this equation.}

\begin{itemize}
\tightlist
\item
  \(T_a\): The author's eight real \(16 \times 16\) Dirac matrices
  \texttt{T16\^{}A{[}a{]}} (\(a = 0, \dots, 7\)), the gamma matrix of
  the frame direction of \(x_a\): \(\gamma^{(x_a)} = T_a\).
\item
  \(I_n\): The \(n \times n\) identity matrix (\(n = 4, 8, 16\); the
  author's \texttt{ID4}, \texttt{ID8}, \texttt{ID16}).
\end{itemize}

\hypertarget{proof-that-the-eight-matrices-are-real}{%
\section{Proof that the eight matrices are
real}\label{proof-that-the-eight-matrices-are-real}}

Every step of the construction (integer entries \(0, \pm1\) of
\(\varepsilon\) and \(\delta\), block assembly, transposition, products
of integer matrices, multiplication by \(\sigma\)) maps integer matrices
to integer matrices, so the \(T_a\) are real (indeed integer) matrices.
The computed entries are only \(-1\), \(0\) and \(+1\)
(\texttt{C1\_real\_integer\_entries} (passed)); each row and each column
of each \(T_a\) holds exactly one non-zero entry, so every \(T_a\) is a
signed permutation matrix (\texttt{C2\_signed\_permutation} (passed)),
hence orthogonal, and (\texttt{C4\_orthogonal\_and\_symmetry} (passed))

\[
T_a^{\mathsf T} = \eta_{aa}\,T_a,\qquad T_a T_a^{\mathsf T} = I_{16}\qquad (a = 0, \dots, 7)
\]

\textbf{Symbols in this equation.}

\begin{itemize}
\tightlist
\item
  \(T_a\): The author's eight real \(16 \times 16\) Dirac matrices
  \texttt{T16\^{}A{[}a{]}} (\(a = 0, \dots, 7\)), the gamma matrix of
  the frame direction of \(x_a\): \(\gamma^{(x_a)} = T_a\).
\item
  \(X^{\mathsf T},\ X^\dagger\): The transpose of a matrix and its
  Hermitian conjugate (transpose of the complex conjugate); for the real
  matrices of this document they coincide.
\item
  \(\eta_{ab}\): The flat frame metric,
  \(\eta = \mathrm{diag}(+1, +1, +1, +1, -1, -1, -1, -1)\) in the order
  \(x_0, \dots, x_7\) (the author's \texttt{eta4488}); \(\eta^{ab}\) is
  the same matrix.
\item
  \(I_n\): The \(n \times n\) identity matrix (\(n = 4, 8, 16\); the
  author's \texttt{ID4}, \texttt{ID8}, \texttt{ID16}).
\end{itemize}

so \(T_0\), \(T_1\), \(T_2\), \(T_3\) are symmetric and \(T_4\),
\(T_5\), \(T_6\), \(T_7\) antisymmetric (the author's
\texttt{Out{[}102{]}}, \texttt{Out{[}103{]}}). Because the matrices are
real, \(T_a^\dagger = T_a^{\mathsf T}\).

\hypertarget{proof-of-the-anticommutation-relations}{%
\section{Proof of the anticommutation
relations}\label{proof-of-the-anticommutation-relations}}

\textbf{Theorem.} For all \(a, b = 0, \dots, 7\)

\[
\{T_a, T_b\} = T_aT_b + T_bT_a = 2\,\eta_{ab}\,I_{16}
\]

\textbf{Symbols in this equation.}

\begin{itemize}
\tightlist
\item
  \(T_a\): The author's eight real \(16 \times 16\) Dirac matrices
  \texttt{T16\^{}A{[}a{]}} (\(a = 0, \dots, 7\)), the gamma matrix of
  the frame direction of \(x_a\): \(\gamma^{(x_a)} = T_a\).
\item
  \([X, Y],\ \{X, Y\}\): The commutator \(XY - YX\) and the
  anticommutator \(XY + YX\) of two matrices.
\item
  \(\eta_{ab}\): The flat frame metric,
  \(\eta = \mathrm{diag}(+1, +1, +1, +1, -1, -1, -1, -1)\) in the order
  \(x_0, \dots, x_7\) (the author's \texttt{eta4488}); \(\eta^{ab}\) is
  the same matrix.
\item
  \(I_n\): The \(n \times n\) identity matrix (\(n = 4, 8, 16\); the
  author's \texttt{ID4}, \texttt{ID8}, \texttt{ID16}).
\end{itemize}

\textbf{Proof.} With the block form of \(T_a\),

\[
T_aT_b + T_bT_a = \begin{pmatrix} \bar\tau_a\tau_b + \bar\tau_b\tau_a & 0 \\ 0 & \tau_a\bar\tau_b + \tau_b\bar\tau_a \end{pmatrix}
\]

\textbf{Symbols in this equation.}

\begin{itemize}
\tightlist
\item
  \(T_a\): The author's eight real \(16 \times 16\) Dirac matrices
  \texttt{T16\^{}A{[}a{]}} (\(a = 0, \dots, 7\)), the gamma matrix of
  the frame direction of \(x_a\): \(\gamma^{(x_a)} = T_a\).
\item
  \(\tau_a\): The author's real \(8 \times 8\) matrices
  \texttt{tau{[}a{]}} (\(a = 0, \dots, 7\)), \(\tau_a = \zeta_a\).
\item
  \(\bar\tau_a\): The author's \texttt{taubar{[}a{]}} (written
  \(\overline{\tau}\) in the notebook), \(\bar\tau_0 = I_8\),
  \(\bar\tau_a = \sigma\,\tau_a^{\mathsf T}\,\sigma\) for
  \(a = 1, \dots, 7\) (not a complex conjugate).
\item
  \(0\): The zero matrix of the size required by its position (in a
  block matrix: a zero block).
\end{itemize}

and the two diagonal blocks are \(2\eta_{ab}I_8\) by the \(8 \times 8\)
relations of the previous section (\texttt{C5\_block\_relations}
(passed)). Those relations are finitely many identities between explicit
integer matrices; they were verified by exact matrix multiplication for
every one of the 64 pairs \((a, b)\), and so were the 64 relations of
the theorem directly (\texttt{C3\_clifford\_relations} (passed)). An
exhaustive exact verification of a finite set of identities between
explicit integer matrices is a complete proof. The author's cell 156
(\texttt{In{[}109{]}}) states the same relations and its stored output
is True in all 64 places (\texttt{B13\_stored\_boolean\_outputs}
(passed)). \(\square\)

In particular \(T_a^2 = +I_{16}\) for the space-like directions
\(x_0, x_1, x_2, x_3\) and \(T_a^2 = -I_{16}\) for the time-like
directions \(x_4, x_5, x_6, x_7\), and different \(T_a\) anticommute.
The \(2^8 = 256\) ordered products \(T_{a_1}\cdots T_{a_k}\)
(\(a_1 < \dots < a_k\), the empty product being \(I_{16}\)) are linearly
independent (\texttt{C30\_basis\_256} (passed)): the real algebra
generated by the \(T_a\) is the whole algebra of real \(16 \times 16\)
matrices, the Clifford algebra \(Cl(4,4) \cong M_{16}(\mathbb R)\); of
the 256 products 136 are symmetric and 120 antisymmetric, as the author
found (\texttt{Out{[}205{]}}, \texttt{Out{[}217{]}},
\texttt{Out{[}220{]}}).

\hypertarget{the-dirac-adjoint-matrix-sigma_16}{%
\section{\texorpdfstring{The Dirac-adjoint matrix
\(\sigma_{16}\)}{The Dirac-adjoint matrix \textbackslash sigma\_\{16\}}}\label{the-dirac-adjoint-matrix-sigma_16}}

The author's \texttt{sigma16} is

\[
\sigma_{16} = T_0T_1T_2T_3 = \begin{pmatrix} -\sigma & 0 \\ 0 & \sigma \end{pmatrix} = \scriptstyle \begin{pmatrix} 0 & 0 & 0 & 0 & -1 & 0 & 0 & 0 & 0 & 0 & 0 & 0 & 0 & 0 & 0 & 0 \\ 0 & 0 & 0 & 0 & 0 & -1 & 0 & 0 & 0 & 0 & 0 & 0 & 0 & 0 & 0 & 0 \\ 0 & 0 & 0 & 0 & 0 & 0 & -1 & 0 & 0 & 0 & 0 & 0 & 0 & 0 & 0 & 0 \\ 0 & 0 & 0 & 0 & 0 & 0 & 0 & -1 & 0 & 0 & 0 & 0 & 0 & 0 & 0 & 0 \\ -1 & 0 & 0 & 0 & 0 & 0 & 0 & 0 & 0 & 0 & 0 & 0 & 0 & 0 & 0 & 0 \\ 0 & -1 & 0 & 0 & 0 & 0 & 0 & 0 & 0 & 0 & 0 & 0 & 0 & 0 & 0 & 0 \\ 0 & 0 & -1 & 0 & 0 & 0 & 0 & 0 & 0 & 0 & 0 & 0 & 0 & 0 & 0 & 0 \\ 0 & 0 & 0 & -1 & 0 & 0 & 0 & 0 & 0 & 0 & 0 & 0 & 0 & 0 & 0 & 0 \\ 0 & 0 & 0 & 0 & 0 & 0 & 0 & 0 & 0 & 0 & 0 & 0 & 1 & 0 & 0 & 0 \\ 0 & 0 & 0 & 0 & 0 & 0 & 0 & 0 & 0 & 0 & 0 & 0 & 0 & 1 & 0 & 0 \\ 0 & 0 & 0 & 0 & 0 & 0 & 0 & 0 & 0 & 0 & 0 & 0 & 0 & 0 & 1 & 0 \\ 0 & 0 & 0 & 0 & 0 & 0 & 0 & 0 & 0 & 0 & 0 & 0 & 0 & 0 & 0 & 1 \\ 0 & 0 & 0 & 0 & 0 & 0 & 0 & 0 & 1 & 0 & 0 & 0 & 0 & 0 & 0 & 0 \\ 0 & 0 & 0 & 0 & 0 & 0 & 0 & 0 & 0 & 1 & 0 & 0 & 0 & 0 & 0 & 0 \\ 0 & 0 & 0 & 0 & 0 & 0 & 0 & 0 & 0 & 0 & 1 & 0 & 0 & 0 & 0 & 0 \\ 0 & 0 & 0 & 0 & 0 & 0 & 0 & 0 & 0 & 0 & 0 & 1 & 0 & 0 & 0 & 0 \end{pmatrix}
\]

\textbf{Symbols in this equation.}

\begin{itemize}
\tightlist
\item
  \(\sigma_{16}\): The author's \texttt{sigma16} \(= T_0 T_1 T_2 T_3\),
  the matrix of the Dirac adjoint
  \(\bar\Psi = \Psi^\dagger\sigma_{16}\).
\item
  \(T_a\): The author's eight real \(16 \times 16\) Dirac matrices
  \texttt{T16\^{}A{[}a{]}} (\(a = 0, \dots, 7\)), the gamma matrix of
  the frame direction of \(x_a\): \(\gamma^{(x_a)} = T_a\).
\item
  \(\sigma\): The real symmetric \(8 \times 8\) matrix
  \(\sigma = \begin{pmatrix} 0 & I_4 \\ I_4 & 0 \end{pmatrix}\) (the
  author's \texttt{sigma88}), with \(\sigma^2 = I_8\).
\item
  \(0\): The zero matrix of the size required by its position (in a
  block matrix: a zero block).
\end{itemize}

It is real, symmetric and its own inverse, and (\texttt{C8\_sigma16}
(passed), \texttt{C9\_sigma16\_conjugation} (passed))

\[
(\sigma_{16}T_a)^{\mathsf T} = -\sigma_{16}T_a,\qquad \sigma_{16}\,T_a\,\sigma_{16}^{-1} = -T_a^{\mathsf T} = -T_a^\dagger\qquad (a = 0, \dots, 7)
\]

\textbf{Symbols in this equation.}

\begin{itemize}
\tightlist
\item
  \(\sigma_{16}\): The author's \texttt{sigma16} \(= T_0 T_1 T_2 T_3\),
  the matrix of the Dirac adjoint
  \(\bar\Psi = \Psi^\dagger\sigma_{16}\).
\item
  \(T_a\): The author's eight real \(16 \times 16\) Dirac matrices
  \texttt{T16\^{}A{[}a{]}} (\(a = 0, \dots, 7\)), the gamma matrix of
  the frame direction of \(x_a\): \(\gamma^{(x_a)} = T_a\).
\item
  \(X^{\mathsf T},\ X^\dagger\): The transpose of a matrix and its
  Hermitian conjugate (transpose of the complex conjugate); for the real
  matrices of this document they coincide.
\end{itemize}

which is the author's check \texttt{In{[}104{]}} (all True). The second
relation is what makes \(\bar\Psi = \Psi^\dagger\sigma_{16}\) the Dirac
adjoint: under \(\Psi \to g\Psi\) with
\(g = \exp(\tfrac12\omega^{ab}S_{ab})\) the bilinear \(\bar\Psi\Psi\) is
invariant and \(\bar\Psi T_a\Psi\) transforms as a vector (because
\(\sigma_{16}S_{ab}\sigma_{16}^{-1} = -S_{ab}^\dagger\) follows from
it).

\[
\bar\Psi = \Psi^\dagger\,\sigma_{16}
\]

\textbf{Symbols in this equation.}

\begin{itemize}
\tightlist
\item
  \(\bar\Psi\): The Dirac adjoint
  \(\bar\Psi = \Psi^\dagger\sigma_{16}\).
\item
  \(\Psi\): The 16-component dirac16complex spinor field (components
  \(\Psi_0, \dots, \Psi_{15}\)), on which the matrices act.
\item
  \(\sigma_{16}\): The author's \texttt{sigma16} \(= T_0 T_1 T_2 T_3\),
  the matrix of the Dirac adjoint
  \(\bar\Psi = \Psi^\dagger\sigma_{16}\).
\item
  \(X^{\mathsf T},\ X^\dagger\): The transpose of a matrix and its
  Hermitian conjugate (transpose of the complex conjugate); for the real
  matrices of this document they coincide.
\end{itemize}

\hypertarget{the-scaled-commutators-spin44-spin44-and-pin44}{%
\section{The scaled commutators: spin(4,4), Spin(4,4) and
Pin(4,4)}\label{the-scaled-commutators-spin44-spin44-and-pin44}}

\hypertarget{definitions}{%
\subsection{Definitions}\label{definitions}}

A \textbf{vector} of the Clifford algebra is \(v = v^aT_a\) with real
\(v^a\); then \(v^2 = \eta_{ab}v^av^b\,I_{16}\), and \(v\) is a
\textbf{unit vector} if \(v^2 = \pm I_{16}\).

\[
v = \sum_{a=0}^{7}v^aT_a,\qquad v^2 = \Big(\sum_{a,b}\eta_{ab}v^av^b\Big)I_{16}
\]

\textbf{Symbols in this equation.}

\begin{itemize}
\tightlist
\item
  \(v,\ v^a\): A vector of the Clifford algebra, \(v = v^a T_a\) with
  real components \(v^a\).
\item
  \(T_a\): The author's eight real \(16 \times 16\) Dirac matrices
  \texttt{T16\^{}A{[}a{]}} (\(a = 0, \dots, 7\)), the gamma matrix of
  the frame direction of \(x_a\): \(\gamma^{(x_a)} = T_a\).
\item
  \(\eta_{ab}\): The flat frame metric,
  \(\eta = \mathrm{diag}(+1, +1, +1, +1, -1, -1, -1, -1)\) in the order
  \(x_0, \dots, x_7\) (the author's \texttt{eta4488}); \(\eta^{ab}\) is
  the same matrix.
\item
  \(I_n\): The \(n \times n\) identity matrix (\(n = 4, 8, 16\); the
  author's \texttt{ID4}, \texttt{ID8}, \texttt{ID16}).
\end{itemize}

\textbf{Pin(4,4)} is the group generated by the unit vectors (products
\(u_1u_2\cdots u_k\) of unit vectors); \textbf{Spin(4,4)} is its
subgroup of products of an even number of unit vectors. The
\textbf{twisted adjoint action}

\[
\rho(g)\,v = \alpha(g)\,v\,g^{-1}
\]

\textbf{Symbols in this equation.}

\begin{itemize}
\tightlist
\item
  \(\rho(g)\): The twisted adjoint action of \(g\) on vectors,
  \(\rho(g)v = \alpha(g)\,v\,g^{-1}\), an element of \(O(4,4)\).
\item
  \(g\): An element of the group \(\mathrm{Pin}(4,4)\) (an invertible
  \(16 \times 16\) real matrix).
\item
  \(\alpha(g)\): The parity automorphism: \(\alpha(g) = g\) for an even
  element (a product of an even number of vectors), \(\alpha(g) = -g\)
  for an odd one.
\item
  \(v,\ v^a\): A vector of the Clifford algebra, \(v = v^a T_a\) with
  real components \(v^a\).
\end{itemize}

maps vectors to vectors and preserves \(v^2\), so
\(\rho(g) \in O(4,4)\); for a unit vector \(u\) it is the reflection in
the hyperplane orthogonal to \(u\), and \(\rho\) is a two-to-one
homomorphism of Pin(4,4) onto \(O(4,4)\) (a standard theorem:
Cartan-Dieudonne together with \(\ker\rho = \{\pm I_{16}\}\)). The Lie
algebra of Pin(4,4), of Spin(4,4) and of their identity component
\(\mathrm{Spin}_0(4,4)\) is the same, \(spin(4,4)\), of dimension
\(\dim O(4,4) = 28\).

\hypertarget{the-scaled-commutators-are-a-basis-of-spin44}{%
\subsection{The scaled commutators are a basis of
spin(4,4)}\label{the-scaled-commutators-are-a-basis-of-spin44}}

\[
S_{ab} = \tfrac14[T_a, T_b] = \tfrac12T_aT_b\quad (a \ne b),\qquad S_{aa} = 0,\qquad S_{ba} = -S_{ab}
\]

\textbf{Symbols in this equation.}

\begin{itemize}
\tightlist
\item
  \(S_{ab}\): The scaled commutator \(S_{ab} = \tfrac14[T_a, T_b]\)
  (\([X, Y] = XY - YX\)), the generators of spin(4,4).
\item
  \(T_a\): The author's eight real \(16 \times 16\) Dirac matrices
  \texttt{T16\^{}A{[}a{]}} (\(a = 0, \dots, 7\)), the gamma matrix of
  the frame direction of \(x_a\): \(\gamma^{(x_a)} = T_a\).
\item
  \([X, Y],\ \{X, Y\}\): The commutator \(XY - YX\) and the
  anticommutator \(XY + YX\) of two matrices.
\item
  \(a, b, c, d\): Frame (orthonormal) indices running over
  \(0, 1, \dots, 7\); the value \(k\) belongs to the author's coordinate
  \(x_k\) (\(x_0\) the hidden space-like direction, \(x_1, x_2, x_3\)
  ordinary space, \(x_4\) the time, \(x_5, x_6, x_7\) the extra
  time-like directions). No sum over a repeated frame index unless a sum
  is written.
\end{itemize}

(\texttt{C22\_SAB\_half\_products} (passed)). The 28 matrices \(S_{ab}\)
with \(a < b\) are linearly independent (\texttt{C21\_SAB\_independent}
(passed)) and satisfy, for all indices
(\texttt{C23\_SAB\_vector\_action} (passed),
\texttt{C24\_spin44\_closure} (passed)),

\[
\begin{aligned}
&[S_{ab}, T_c] = \eta_{bc}T_a - \eta_{ac}T_b,\\
&[S_{ab}, S_{cd}] = \eta_{bc}S_{ad} - \eta_{ac}S_{bd} - \eta_{bd}S_{ac} + \eta_{ad}S_{bc}
\end{aligned}
\]

\textbf{Symbols in this equation.}

\begin{itemize}
\tightlist
\item
  \(S_{ab}\): The scaled commutator \(S_{ab} = \tfrac14[T_a, T_b]\)
  (\([X, Y] = XY - YX\)), the generators of spin(4,4).
\item
  \(T_a\): The author's eight real \(16 \times 16\) Dirac matrices
  \texttt{T16\^{}A{[}a{]}} (\(a = 0, \dots, 7\)), the gamma matrix of
  the frame direction of \(x_a\): \(\gamma^{(x_a)} = T_a\).
\item
  \([X, Y],\ \{X, Y\}\): The commutator \(XY - YX\) and the
  anticommutator \(XY + YX\) of two matrices.
\item
  \(\eta_{ab}\): The flat frame metric,
  \(\eta = \mathrm{diag}(+1, +1, +1, +1, -1, -1, -1, -1)\) in the order
  \(x_0, \dots, x_7\) (the author's \texttt{eta4488}); \(\eta^{ab}\) is
  the same matrix.
\item
  \(a, b, c, d\): Frame (orthonormal) indices running over
  \(0, 1, \dots, 7\); the value \(k\) belongs to the author's coordinate
  \(x_k\) (\(x_0\) the hidden space-like direction, \(x_1, x_2, x_3\)
  ordinary space, \(x_4\) the time, \(x_5, x_6, x_7\) the extra
  time-like directions). No sum over a repeated frame index unless a sum
  is written.
\end{itemize}

The first relation follows from the theorem of the previous section in
one line:
\([T_aT_b, T_c] = T_a\{T_b, T_c\} - \{T_a, T_c\}T_b = 2\eta_{bc}T_a - 2\eta_{ac}T_b\).
It says that \(\mathrm{ad}(S_{ab})\) acts on the span of the \(T_c\) by
the matrix

\[
(L_{ab})^c{}_d = \delta^c_a\,\eta_{bd} - \delta^c_b\,\eta_{ad},\qquad L_{ab}^{\mathsf T}\eta + \eta L_{ab} = 0
\]

\textbf{Symbols in this equation.}

\begin{itemize}
\tightlist
\item
  \(L_{ab}\): The generator of \(O(4,4)\) on vectors in the plane
  \((a, b)\):
  \((L_{ab})^c{}_d = \delta^c_a\eta_{bd} - \delta^c_b\eta_{ad}\).
\item
  \(\delta_{pq}\): The Kronecker delta: \(1\) if \(p = q\), else \(0\)
  (the author's \texttt{ID4{[}{[}p,q{]}{]}}).
\item
  \(\eta_{ab}\): The flat frame metric,
  \(\eta = \mathrm{diag}(+1, +1, +1, +1, -1, -1, -1, -1)\) in the order
  \(x_0, \dots, x_7\) (the author's \texttt{eta4488}); \(\eta^{ab}\) is
  the same matrix.
\item
  \(a, b, c, d\): Frame (orthonormal) indices running over
  \(0, 1, \dots, 7\); the value \(k\) belongs to the author's coordinate
  \(x_k\) (\(x_0\) the hidden space-like direction, \(x_1, x_2, x_3\)
  ordinary space, \(x_4\) the time, \(x_5, x_6, x_7\) the extra
  time-like directions). No sum over a repeated frame index unless a sum
  is written.
\end{itemize}

the generator of \(O(4,4)\) in the plane \((a, b)\)
(\texttt{C25\_isomorphism\_so44} (passed)). The second relation, with
the structure constants of \(so(4,4)\), shows that
\(S_{ab} \mapsto L_{ab}\) is an isomorphism of Lie algebras
\(\mathrm{span}\{S_{ab}\} \cong so(4,4)\). Since every
\(\exp(\theta S_{ab})\) lies in Pin(4,4) (next paragraph) and
\(\tfrac{\mathrm{d}}{\mathrm{d}\theta}\exp(\theta S_{ab})|_0 = S_{ab}\),
the 28 independent \(S_{ab}\) lie in the 28-dimensional Lie algebra of
Pin(4,4) and are a basis of it:
\(spin(4,4) = \mathrm{span}_{\mathbb R}\{S_{ab}\}\).

\hypertarget{the-exponentials-are-products-of-two-unit-vectors}{%
\subsection{The exponentials are products of two unit
vectors}\label{the-exponentials-are-products-of-two-unit-vectors}}

\((T_aT_b)^2 = -T_a^2T_b^2 = -\eta_{aa}\eta_{bb}I_{16}\) for
\(a \ne b\), so the exponential series sums to
(\texttt{C26\_exponentials\_two\_unit\_vectors} (passed))

\[
\exp(\theta S_{ab}) = \begin{cases} \cos\tfrac\theta2\,I_{16} + \sin\tfrac\theta2\,T_aT_b, & \eta_{aa} = \eta_{bb}\ \text{(a rotation)},\\ \cosh\tfrac\theta2\,I_{16} + \sinh\tfrac\theta2\,T_aT_b, & \eta_{aa} = -\eta_{bb}\ \text{(a boost)} \end{cases}
\]

\textbf{Symbols in this equation.}

\begin{itemize}
\tightlist
\item
  \(S_{ab}\): The scaled commutator \(S_{ab} = \tfrac14[T_a, T_b]\)
  (\([X, Y] = XY - YX\)), the generators of spin(4,4).
\item
  \(T_a\): The author's eight real \(16 \times 16\) Dirac matrices
  \texttt{T16\^{}A{[}a{]}} (\(a = 0, \dots, 7\)), the gamma matrix of
  the frame direction of \(x_a\): \(\gamma^{(x_a)} = T_a\).
\item
  \(\theta\): A real parameter (an angle for a rotation, a rapidity for
  a boost).
\item
  \(\eta_{ab}\): The flat frame metric,
  \(\eta = \mathrm{diag}(+1, +1, +1, +1, -1, -1, -1, -1)\) in the order
  \(x_0, \dots, x_7\) (the author's \texttt{eta4488}); \(\eta^{ab}\) is
  the same matrix.
\item
  \(I_n\): The \(n \times n\) identity matrix (\(n = 4, 8, 16\); the
  author's \texttt{ID4}, \texttt{ID8}, \texttt{ID16}).
\end{itemize}

and each is a product \(u_1u_2\) of two unit vectors: for
\(\eta_{aa} = +1\), \(u_1 = T_a\) and
\(u_2 = \cos\tfrac\theta2T_a + \sin\tfrac\theta2T_b\) (rotation) or
\(u_2 = \cosh\tfrac\theta2T_a + \sinh\tfrac\theta2T_b\) (boost); for
\(\eta_{aa} = -1\), \(u_1 = -T_a\) and the sign of the \(T_b\) term
reversed. Hence
\(\exp(\theta S_{ab}) \in \mathrm{Spin}(4,4) \subset \mathrm{Pin}(4,4)\).

\hypertarget{what-the-exponentials-generate}{%
\subsection{What the exponentials
generate}\label{what-the-exponentials-generate}}

\textbf{Theorem.} The group generated by
\(\{\exp(\theta S_{ab}) : \theta \in \mathbb R,\ 0 \le a < b \le 7\}\)
is the identity component \(\mathrm{Spin}_0(4,4)\) of Pin(4,4).

\textbf{Proof.} (i) Every \(\exp(\theta S_{ab})\) is joined to
\(I_{16}\) by the path \(\theta' \mapsto \exp(\theta' S_{ab})\),
\(0 \le \theta' \le \theta\), inside Pin(4,4), so it and every product
of such elements lie in the identity component; their images
\(\rho = \exp(\theta L_{ab})\) in \(O(4,4)\) have the determinants of
both diagonal blocks (the space-like block \(x_0..x_3\) and the
time-like block \(x_4..x_7\)) equal to \(+1\) for rotations and
\(\cosh\theta > 0\) for boosts, i.e.~they lie in the identity component
\(SO_0(4,4)\) (\texttt{C28\_exponentials\_in\_identity\_component}
(passed)). (ii) The map
\((\theta_1, \dots, \theta_{28}) \mapsto \prod_k\exp(\theta_kS_k)\),
\(S_k\) running over the 28 basis elements, has at \(0\) the
differential \((\theta_k) \mapsto \sum_k\theta_kS_k\), an isomorphism
onto the Lie algebra; by the inverse function theorem its image contains
a neighbourhood of \(I_{16}\) in \(\mathrm{Spin}_0(4,4)\). A connected
group is generated by any neighbourhood of its identity, so the
generated group is all of \(\mathrm{Spin}_0(4,4)\). \(\square\)

\textbf{The four components of Pin(4,4).} The twisted adjoint images of
\(I_{16}\), \(T_0\), \(T_4\) and \(T_0T_4\) in \(O(4,4)\) have the
following signs of the determinants of the space-like block
(\(x_0..x_3\)) and of the time-like block (\(x_4..x_7\))
(\texttt{C27\_pin44\_components} (passed)):

\begin{longtable}[]{@{}
  >{\raggedright\arraybackslash}p{(\columnwidth - 8\tabcolsep) * \real{0.2000}}
  >{\raggedright\arraybackslash}p{(\columnwidth - 8\tabcolsep) * \real{0.2000}}
  >{\raggedright\arraybackslash}p{(\columnwidth - 8\tabcolsep) * \real{0.2000}}
  >{\raggedright\arraybackslash}p{(\columnwidth - 8\tabcolsep) * \real{0.2000}}
  >{\raggedright\arraybackslash}p{(\columnwidth - 8\tabcolsep) * \real{0.2000}}@{}}
\toprule\noalign{}
\begin{minipage}[b]{\linewidth}\raggedright
element \(g\)
\end{minipage} & \begin{minipage}[b]{\linewidth}\raggedright
sign det (space-like block)
\end{minipage} & \begin{minipage}[b]{\linewidth}\raggedright
sign det (time-like block)
\end{minipage} & \begin{minipage}[b]{\linewidth}\raggedright
det \(\rho(g)\)
\end{minipage} & \begin{minipage}[b]{\linewidth}\raggedright
component
\end{minipage} \\
\midrule\noalign{}
\endhead
\bottomrule\noalign{}
\endlastfoot
\(I_{16}\) & +1 & +1 & +1 & \(\mathrm{Spin}_0(4,4)\) \\
\(T_0\) & -1 & +1 & -1 & \(\mathrm{Spin}_0(4,4)\,T_0\) \\
\(T_4\) & +1 & -1 & -1 & \(\mathrm{Spin}_0(4,4)\,T_4\) \\
\(T_0T_4\) & -1 & -1 & +1 & \(\mathrm{Spin}_0(4,4)\,T_0T_4\) \\
\end{longtable}

The four sign pairs are different, and the block signs are constant on
connected pieces of \(O(4,4)\), so the four elements lie in the four
different components of Pin(4,4), each of which is a coset of the
identity component:

\[
\mathrm{Pin}(4,4) = \mathrm{Spin}_0(4,4)\ \cup\ \mathrm{Spin}_0(4,4)\,T_0\ \cup\ \mathrm{Spin}_0(4,4)\,T_4\ \cup\ \mathrm{Spin}_0(4,4)\,T_0T_4 = \big\langle \exp(\theta S_{ab}),\ T_0,\ T_4\big\rangle
\]

\textbf{Symbols in this equation.}

\begin{itemize}
\tightlist
\item
  \(S_{ab}\): The scaled commutator \(S_{ab} = \tfrac14[T_a, T_b]\)
  (\([X, Y] = XY - YX\)), the generators of spin(4,4).
\item
  \(T_a\): The author's eight real \(16 \times 16\) Dirac matrices
  \texttt{T16\^{}A{[}a{]}} (\(a = 0, \dots, 7\)), the gamma matrix of
  the frame direction of \(x_a\): \(\gamma^{(x_a)} = T_a\).
\item
  \(\theta\): A real parameter (an angle for a rotation, a rapidity for
  a boost).
\end{itemize}

(\(\langle\cdots\rangle\) is the group generated). Spin(4,4) is the
union of the first and the last coset. So the scaled commutators
generate Pin(4,4) as its Lie algebra and, through their exponentials,
its identity component; the reflections \(T_0\) (space-like) and \(T_4\)
(time-like) supply the other three components, which no product of
exponentials can reach.

\textbf{Irreducibility.} The only matrices commuting with all eight
\(T_a\) are the multiples of \(I_{16}\)
(\texttt{C18\_pin44\_irreducible} (passed)): \(\mathbb C^{16}\) is an
irreducible representation of Pin(4,4). The matrices commuting with all
\(S_{ab}\) form the two-dimensional algebra spanned by \(I_{16}\) and
\(\Gamma\) (\texttt{C17\_commutant\_spin44} (passed)); under Spin(4,4)
the space splits into the two halves \(P_L\mathbb C^{16}\) and
\(P_R\mathbb C^{16}\), each irreducible, and the two are inequivalent
(no non-zero intertwiner;
\texttt{C19\_halves\_irreducible\_inequivalent} (passed)).

\hypertarget{chirality-re-derivation-of-the-projection-operators}{%
\section{Chirality: re-derivation of the projection
operators}\label{chirality-re-derivation-of-the-projection-operators}}

The author's subsection ``Left - Handed vs . Right - Handed Fermion
Fields (Chirality) ; projection operators'' (cells 160 to 166) reads,
verbatim:

\begin{verbatim}
cell 160  Subsection 
    Left - Handed vs . Right - Handed Fermion Fields (Chirality) ; projection operators
cell 161  Input In[111]:=
    (Subscript[P, L]=1/2(ID16-T16^A[8]))//MatrixForm
cell 162  Output Out[111]//MatrixForm=
    (\[NoBreak]{{1, 0, 0, 0, 0, 0, 0, 0, 0, 0, 0, 0, 0, 0, 0, 0}, {0, 1, 0, 0, 0, 0, 0, 0, 0, 0, 0, 0, 0, 0, 0, 0}, {0, 0, 1 ...
cell 163  Input In[112]:=
    (Subscript[P, R]=1/2(ID16+T16^A[8]))//MatrixForm
cell 164  Output Out[112]//MatrixForm=
    (\[NoBreak]{{0, 0, 0, 0, 0, 0, 0, 0, 0, 0, 0, 0, 0, 0, 0, 0}, {0, 0, 0, 0, 0, 0, 0, 0, 0, 0, 0, 0, 0, 0, 0, 0}, {0, 0, 0 ...
cell 165  Input In[113]:=
    {ID16== Subscript[P, L]+Subscript[P, R],Subscript[P, L].Subscript[P, L]==Subscript[P, L],Subscript[P, L].Subscript[P, R]==Subscript[P, R].Subscript[P, L]==ZERO16}
cell 166  Output Out[113]=
    {True,True,True} ...
\end{verbatim}

\textbf{Step 1: the chirality matrix.} A chirality operator must commute
with every Lorentz generator \(S_{ab}\) (so that it is the same in every
frame), square to the identity, and not be a multiple of \(I_{16}\). The
commutant of the \(S_{ab}\) is \(\{xI_{16} + y\Gamma\}\) with
(\texttt{C17\_commutant\_spin44} (passed),
\texttt{C10\_chirality\_matrix} (passed),
\texttt{C11\_chirality\_commutes\_with\_SAB} (passed))

\[
\Gamma = T_0T_1T_2T_3T_4T_5T_6T_7,\qquad \Gamma^2 = (-1)^{8\cdot7/2}\,\eta_{00}\eta_{11}\cdots\eta_{77}\,I_{16} = I_{16},\qquad \Gamma T_a = -T_a\Gamma
\]

\textbf{Symbols in this equation.}

\begin{itemize}
\tightlist
\item
  \(\Gamma\): The chirality matrix, the author's
  \texttt{T16\^{}A{[}8{]}} \(= T_0 T_1 \cdots T_7\).
\item
  \(T_a\): The author's eight real \(16 \times 16\) Dirac matrices
  \texttt{T16\^{}A{[}a{]}} (\(a = 0, \dots, 7\)), the gamma matrix of
  the frame direction of \(x_a\): \(\gamma^{(x_a)} = T_a\).
\item
  \(\eta_{ab}\): The flat frame metric,
  \(\eta = \mathrm{diag}(+1, +1, +1, +1, -1, -1, -1, -1)\) in the order
  \(x_0, \dots, x_7\) (the author's \texttt{eta4488}); \(\eta^{ab}\) is
  the same matrix.
\item
  \(I_n\): The \(n \times n\) identity matrix (\(n = 4, 8, 16\); the
  author's \texttt{ID4}, \texttt{ID8}, \texttt{ID16}).
\end{itemize}

(moving the eight factors of \(\Gamma\Gamma\) past each other gives the
sign \((-1)^{28} = +1\), and the product of the squares is
\((+1)^4(-1)^4 = +1\); \texttt{C12\_chirality\_square\_formula}
(passed)). \(\Gamma\) anticommutes with every \(T_a\) because \(T_a\)
anticommutes with the seven other factors and commutes with itself
(seven sign changes). Explicitly (\texttt{B8\_stored\_T16A8} (passed))

\[
\Gamma = \begin{pmatrix} -I_8 & 0 \\ 0 & I_8 \end{pmatrix} = \scriptstyle \begin{pmatrix} -1 & 0 & 0 & 0 & 0 & 0 & 0 & 0 & 0 & 0 & 0 & 0 & 0 & 0 & 0 & 0 \\ 0 & -1 & 0 & 0 & 0 & 0 & 0 & 0 & 0 & 0 & 0 & 0 & 0 & 0 & 0 & 0 \\ 0 & 0 & -1 & 0 & 0 & 0 & 0 & 0 & 0 & 0 & 0 & 0 & 0 & 0 & 0 & 0 \\ 0 & 0 & 0 & -1 & 0 & 0 & 0 & 0 & 0 & 0 & 0 & 0 & 0 & 0 & 0 & 0 \\ 0 & 0 & 0 & 0 & -1 & 0 & 0 & 0 & 0 & 0 & 0 & 0 & 0 & 0 & 0 & 0 \\ 0 & 0 & 0 & 0 & 0 & -1 & 0 & 0 & 0 & 0 & 0 & 0 & 0 & 0 & 0 & 0 \\ 0 & 0 & 0 & 0 & 0 & 0 & -1 & 0 & 0 & 0 & 0 & 0 & 0 & 0 & 0 & 0 \\ 0 & 0 & 0 & 0 & 0 & 0 & 0 & -1 & 0 & 0 & 0 & 0 & 0 & 0 & 0 & 0 \\ 0 & 0 & 0 & 0 & 0 & 0 & 0 & 0 & 1 & 0 & 0 & 0 & 0 & 0 & 0 & 0 \\ 0 & 0 & 0 & 0 & 0 & 0 & 0 & 0 & 0 & 1 & 0 & 0 & 0 & 0 & 0 & 0 \\ 0 & 0 & 0 & 0 & 0 & 0 & 0 & 0 & 0 & 0 & 1 & 0 & 0 & 0 & 0 & 0 \\ 0 & 0 & 0 & 0 & 0 & 0 & 0 & 0 & 0 & 0 & 0 & 1 & 0 & 0 & 0 & 0 \\ 0 & 0 & 0 & 0 & 0 & 0 & 0 & 0 & 0 & 0 & 0 & 0 & 1 & 0 & 0 & 0 \\ 0 & 0 & 0 & 0 & 0 & 0 & 0 & 0 & 0 & 0 & 0 & 0 & 0 & 1 & 0 & 0 \\ 0 & 0 & 0 & 0 & 0 & 0 & 0 & 0 & 0 & 0 & 0 & 0 & 0 & 0 & 1 & 0 \\ 0 & 0 & 0 & 0 & 0 & 0 & 0 & 0 & 0 & 0 & 0 & 0 & 0 & 0 & 0 & 1 \end{pmatrix}
\]

\textbf{Symbols in this equation.}

\begin{itemize}
\tightlist
\item
  \(\Gamma\): The chirality matrix, the author's
  \texttt{T16\^{}A{[}8{]}} \(= T_0 T_1 \cdots T_7\).
\item
  \(I_n\): The \(n \times n\) identity matrix (\(n = 4, 8, 16\); the
  author's \texttt{ID4}, \texttt{ID8}, \texttt{ID16}).
\item
  \(0\): The zero matrix of the size required by its position (in a
  block matrix: a zero block).
\end{itemize}

\textbf{Step 2: the projectors.} A projector \(X = xI_{16} + y\Gamma\)
satisfies \(X^2 = X\), i.e.~\(x^2 + y^2 = x\) and \(2xy = y\); the
solutions are \((x, y) = (0, 0)\), \((1, 0)\), \((\tfrac12, -\tfrac12)\)
and \((\tfrac12, \tfrac12)\)
(\texttt{C20\_idempotents\_of\_chirality\_algebra} (passed)). The two
non-trivial ones are the author's

\[
P_L = \tfrac12\big(I_{16} - \Gamma\big) = \begin{pmatrix} I_8 & 0 \\ 0 & 0 \end{pmatrix},\qquad P_R = \tfrac12\big(I_{16} + \Gamma\big) = \begin{pmatrix} 0 & 0 \\ 0 & I_8 \end{pmatrix}
\]

\textbf{Symbols in this equation.}

\begin{itemize}
\tightlist
\item
  \(P_L,\ P_R\): The author's chiral projection operators \texttt{P\_L}
  \(= \tfrac12(I_{16} - \Gamma)\) and \texttt{P\_R}
  \(= \tfrac12(I_{16} + \Gamma)\) (left-handed and right-handed).
\item
  \(\Gamma\): The chirality matrix, the author's
  \texttt{T16\^{}A{[}8{]}} \(= T_0 T_1 \cdots T_7\).
\item
  \(I_n\): The \(n \times n\) identity matrix (\(n = 4, 8, 16\); the
  author's \texttt{ID4}, \texttt{ID8}, \texttt{ID16}).
\item
  \(0\): The zero matrix of the size required by its position (in a
  block matrix: a zero block).
\end{itemize}

equal, entry by entry, to the matrices the notebook stored
(\texttt{B9\_stored\_PL} (passed), \texttt{B10\_stored\_PR} (passed)):

\[
P_L = \scriptstyle \begin{pmatrix} 1 & 0 & 0 & 0 & 0 & 0 & 0 & 0 & 0 & 0 & 0 & 0 & 0 & 0 & 0 & 0 \\ 0 & 1 & 0 & 0 & 0 & 0 & 0 & 0 & 0 & 0 & 0 & 0 & 0 & 0 & 0 & 0 \\ 0 & 0 & 1 & 0 & 0 & 0 & 0 & 0 & 0 & 0 & 0 & 0 & 0 & 0 & 0 & 0 \\ 0 & 0 & 0 & 1 & 0 & 0 & 0 & 0 & 0 & 0 & 0 & 0 & 0 & 0 & 0 & 0 \\ 0 & 0 & 0 & 0 & 1 & 0 & 0 & 0 & 0 & 0 & 0 & 0 & 0 & 0 & 0 & 0 \\ 0 & 0 & 0 & 0 & 0 & 1 & 0 & 0 & 0 & 0 & 0 & 0 & 0 & 0 & 0 & 0 \\ 0 & 0 & 0 & 0 & 0 & 0 & 1 & 0 & 0 & 0 & 0 & 0 & 0 & 0 & 0 & 0 \\ 0 & 0 & 0 & 0 & 0 & 0 & 0 & 1 & 0 & 0 & 0 & 0 & 0 & 0 & 0 & 0 \\ 0 & 0 & 0 & 0 & 0 & 0 & 0 & 0 & 0 & 0 & 0 & 0 & 0 & 0 & 0 & 0 \\ 0 & 0 & 0 & 0 & 0 & 0 & 0 & 0 & 0 & 0 & 0 & 0 & 0 & 0 & 0 & 0 \\ 0 & 0 & 0 & 0 & 0 & 0 & 0 & 0 & 0 & 0 & 0 & 0 & 0 & 0 & 0 & 0 \\ 0 & 0 & 0 & 0 & 0 & 0 & 0 & 0 & 0 & 0 & 0 & 0 & 0 & 0 & 0 & 0 \\ 0 & 0 & 0 & 0 & 0 & 0 & 0 & 0 & 0 & 0 & 0 & 0 & 0 & 0 & 0 & 0 \\ 0 & 0 & 0 & 0 & 0 & 0 & 0 & 0 & 0 & 0 & 0 & 0 & 0 & 0 & 0 & 0 \\ 0 & 0 & 0 & 0 & 0 & 0 & 0 & 0 & 0 & 0 & 0 & 0 & 0 & 0 & 0 & 0 \\ 0 & 0 & 0 & 0 & 0 & 0 & 0 & 0 & 0 & 0 & 0 & 0 & 0 & 0 & 0 & 0 \end{pmatrix}
\]

\[
P_R = \scriptstyle \begin{pmatrix} 0 & 0 & 0 & 0 & 0 & 0 & 0 & 0 & 0 & 0 & 0 & 0 & 0 & 0 & 0 & 0 \\ 0 & 0 & 0 & 0 & 0 & 0 & 0 & 0 & 0 & 0 & 0 & 0 & 0 & 0 & 0 & 0 \\ 0 & 0 & 0 & 0 & 0 & 0 & 0 & 0 & 0 & 0 & 0 & 0 & 0 & 0 & 0 & 0 \\ 0 & 0 & 0 & 0 & 0 & 0 & 0 & 0 & 0 & 0 & 0 & 0 & 0 & 0 & 0 & 0 \\ 0 & 0 & 0 & 0 & 0 & 0 & 0 & 0 & 0 & 0 & 0 & 0 & 0 & 0 & 0 & 0 \\ 0 & 0 & 0 & 0 & 0 & 0 & 0 & 0 & 0 & 0 & 0 & 0 & 0 & 0 & 0 & 0 \\ 0 & 0 & 0 & 0 & 0 & 0 & 0 & 0 & 0 & 0 & 0 & 0 & 0 & 0 & 0 & 0 \\ 0 & 0 & 0 & 0 & 0 & 0 & 0 & 0 & 0 & 0 & 0 & 0 & 0 & 0 & 0 & 0 \\ 0 & 0 & 0 & 0 & 0 & 0 & 0 & 0 & 1 & 0 & 0 & 0 & 0 & 0 & 0 & 0 \\ 0 & 0 & 0 & 0 & 0 & 0 & 0 & 0 & 0 & 1 & 0 & 0 & 0 & 0 & 0 & 0 \\ 0 & 0 & 0 & 0 & 0 & 0 & 0 & 0 & 0 & 0 & 1 & 0 & 0 & 0 & 0 & 0 \\ 0 & 0 & 0 & 0 & 0 & 0 & 0 & 0 & 0 & 0 & 0 & 1 & 0 & 0 & 0 & 0 \\ 0 & 0 & 0 & 0 & 0 & 0 & 0 & 0 & 0 & 0 & 0 & 0 & 1 & 0 & 0 & 0 \\ 0 & 0 & 0 & 0 & 0 & 0 & 0 & 0 & 0 & 0 & 0 & 0 & 0 & 1 & 0 & 0 \\ 0 & 0 & 0 & 0 & 0 & 0 & 0 & 0 & 0 & 0 & 0 & 0 & 0 & 0 & 1 & 0 \\ 0 & 0 & 0 & 0 & 0 & 0 & 0 & 0 & 0 & 0 & 0 & 0 & 0 & 0 & 0 & 1 \end{pmatrix}
\]

\textbf{Symbols in this equation.}

\begin{itemize}
\tightlist
\item
  \(P_L,\ P_R\): The author's chiral projection operators \texttt{P\_L}
  \(= \tfrac12(I_{16} - \Gamma)\) and \texttt{P\_R}
  \(= \tfrac12(I_{16} + \Gamma)\) (left-handed and right-handed).
\end{itemize}

\textbf{Step 3: their properties} (\texttt{C13\_projectors} (passed),
\texttt{C14\_projectors\_explicit} (passed),
\texttt{C15\_projectors\_and\_gammas} (passed),
\texttt{C16\_projectors\_and\_adjoint} (passed)): the author's three
tests of cell 165 and more,

\[
\begin{aligned}
&P_L + P_R = I_{16},\quad P_L^2 = P_L,\quad P_R^2 = P_R,\quad P_LP_R = P_RP_L = 0,\quad \mathrm{Tr}\,P_L = \mathrm{Tr}\,P_R = 8,\\
&P_LT_a = T_aP_R,\qquad [P_L, S_{ab}] = 0,\qquad [P_L, \sigma_{16}] = 0,\qquad P_L^{\mathsf T} = P_L
\end{aligned}
\]

\textbf{Symbols in this equation.}

\begin{itemize}
\tightlist
\item
  \(P_L,\ P_R\): The author's chiral projection operators \texttt{P\_L}
  \(= \tfrac12(I_{16} - \Gamma)\) and \texttt{P\_R}
  \(= \tfrac12(I_{16} + \Gamma)\) (left-handed and right-handed).
\item
  \(T_a\): The author's eight real \(16 \times 16\) Dirac matrices
  \texttt{T16\^{}A{[}a{]}} (\(a = 0, \dots, 7\)), the gamma matrix of
  the frame direction of \(x_a\): \(\gamma^{(x_a)} = T_a\).
\item
  \(S_{ab}\): The scaled commutator \(S_{ab} = \tfrac14[T_a, T_b]\)
  (\([X, Y] = XY - YX\)), the generators of spin(4,4).
\item
  \(\sigma_{16}\): The author's \texttt{sigma16} \(= T_0 T_1 T_2 T_3\),
  the matrix of the Dirac adjoint
  \(\bar\Psi = \Psi^\dagger\sigma_{16}\).
\item
  \([X, Y],\ \{X, Y\}\): The commutator \(XY - YX\) and the
  anticommutator \(XY + YX\) of two matrices.
\item
  \(X^{\mathsf T},\ X^\dagger\): The transpose of a matrix and its
  Hermitian conjugate (transpose of the complex conjugate); for the real
  matrices of this document they coincide.
\item
  \(I_n\): The \(n \times n\) identity matrix (\(n = 4, 8, 16\); the
  author's \texttt{ID4}, \texttt{ID8}, \texttt{ID16}).
\end{itemize}

So the left-handed part \(P_L\Psi\) consists of the components
\(\Psi_0, \dots, \Psi_7\) and the right-handed part \(P_R\Psi\) of
\(\Psi_8, \dots, \Psi_{15}\); each is an irreducible representation of
Spin(4,4), the two inequivalent; every Dirac matrix exchanges them. Two
consequences matter for every Lagrangian built with these matrices:
because \(\sigma_{16}\) commutes with \(P_L\),

\[
\overline{P_L\Psi} = \bar\Psi P_L,\qquad \bar\Psi\Psi = \bar\Psi P_L\Psi + \bar\Psi P_R\Psi,\qquad \bar\Psi T_a\Psi = \bar\Psi P_LT_aP_R\Psi + \bar\Psi P_RT_aP_L\Psi
\]

\textbf{Symbols in this equation.}

\begin{itemize}
\tightlist
\item
  \(\bar\Psi\): The Dirac adjoint
  \(\bar\Psi = \Psi^\dagger\sigma_{16}\).
\item
  \(\Psi\): The 16-component dirac16complex spinor field (components
  \(\Psi_0, \dots, \Psi_{15}\)), on which the matrices act.
\item
  \(P_L,\ P_R\): The author's chiral projection operators \texttt{P\_L}
  \(= \tfrac12(I_{16} - \Gamma)\) and \texttt{P\_R}
  \(= \tfrac12(I_{16} + \Gamma)\) (left-handed and right-handed).
\item
  \(T_a\): The author's eight real \(16 \times 16\) Dirac matrices
  \texttt{T16\^{}A{[}a{]}} (\(a = 0, \dots, 7\)), the gamma matrix of
  the frame direction of \(x_a\): \(\gamma^{(x_a)} = T_a\).
\end{itemize}

a bilinear with an even number of Dirac matrices (such as the mass term
\(\bar\Psi\Psi\)) never mixes the two chiralities, while one with an odd
number (such as the kinetic term \(\bar\Psi T_a\partial\Psi\)) always
connects the left-handed with the right-handed part. In this signature
the Dirac adjoint is built with the even matrix \(\sigma_{16}\); in
four-dimensional Minkowski space it is built with the odd \(\gamma^0\),
which is why there the mass term mixes the chiralities and the kinetic
term does not.

\hypertarget{other-important-16-times-16-matrices-of-the-notebook}{%
\section{\texorpdfstring{Other important \(16 \times 16\) matrices of
the
notebook}{Other important 16 \textbackslash times 16 matrices of the notebook}}\label{other-important-16-times-16-matrices-of-the-notebook}}

\textbf{The complex structure} \(J = T_4\) (the author's cells 654 to
669): \(J^2 = -I_{16}\) and \(J^{\mathsf T}\sigma_{16}J = \sigma_{16}\),
the author's \texttt{In{[}476{]}} (\texttt{C29\_complex\_structure\_T4}
(passed)); the author notes that he uses \(x_4\) (not the periodic
\(x_0\)) for it.

\[
J = T_4,\qquad J^2 = -I_{16},\qquad J^{\mathsf T}\sigma_{16}J = \sigma_{16}
\]

\textbf{Symbols in this equation.}

\begin{itemize}
\tightlist
\item
  \(J\): The author's complex structure \(J = T_4\) (the Dirac matrix of
  the time \(x_4\)), \(J^2 = -I_{16}\).
\item
  \(T_a\): The author's eight real \(16 \times 16\) Dirac matrices
  \texttt{T16\^{}A{[}a{]}} (\(a = 0, \dots, 7\)), the gamma matrix of
  the frame direction of \(x_a\): \(\gamma^{(x_a)} = T_a\).
\item
  \(\sigma_{16}\): The author's \texttt{sigma16} \(= T_0 T_1 T_2 T_3\),
  the matrix of the Dirac adjoint
  \(\bar\Psi = \Psi^\dagger\sigma_{16}\).
\item
  \(X^{\mathsf T},\ X^\dagger\): The transpose of a matrix and its
  Hermitian conjugate (transpose of the complex conjugate); for the real
  matrices of this document they coincide.
\item
  \(I_n\): The \(n \times n\) identity matrix (\(n = 4, 8, 16\); the
  author's \texttt{ID4}, \texttt{ID8}, \texttt{ID16}).
\end{itemize}

\textbf{The product} \(T_1T_2T_3T_4T_5T_6\) stored by the author
(\texttt{Out{[}489{]}}) is reproduced (\texttt{B12\_stored\_T1\_to\_T6}
(passed)); the products \(T_4T_5T_6T_7T_a\) are symmetric
(\texttt{Out{[}100{]}}) and \(\Gamma = \sigma_{16}T_4T_5T_6T_7\)
(\texttt{Out{[}101{]}}, \texttt{Out{[}498{]}}), all reproduced
(\texttt{B13\_stored\_boolean\_outputs} (passed)).

\textbf{The basis of all real \(16 \times 16\) matrices} is the set of
the 256 ordered products of the \(T_a\), with
\(\mathrm{Tr}(MM)/16 = \pm1\) for each (the author's
\texttt{Out{[}202{]}}) and 136 symmetric, 120 antisymmetric members
(\texttt{C30\_basis\_256} (passed)).

\hypertarget{agreement-with-every-stored-output-of-the-matrix-cells}{%
\section{Agreement with every stored output of the matrix
cells}\label{agreement-with-every-stored-output-of-the-matrix-cells}}

The notebook file keeps the outputs that Mathematica computed when the
author evaluated the cells. Each was parsed and compared with the
rebuilt matrices (matrices entry by entry; True/False lists value by
value):

\begin{longtable}[]{@{}
  >{\raggedright\arraybackslash}p{(\columnwidth - 6\tabcolsep) * \real{0.0600}}
  >{\raggedright\arraybackslash}p{(\columnwidth - 6\tabcolsep) * \real{0.2400}}
  >{\raggedright\arraybackslash}p{(\columnwidth - 6\tabcolsep) * \real{0.5800}}
  >{\raggedright\arraybackslash}p{(\columnwidth - 6\tabcolsep) * \real{0.1200}}@{}}
\toprule\noalign{}
\begin{minipage}[b]{\linewidth}\raggedright
cell
\end{minipage} & \begin{minipage}[b]{\linewidth}\raggedright
label
\end{minipage} & \begin{minipage}[b]{\linewidth}\raggedright
what
\end{minipage} & \begin{minipage}[b]{\linewidth}\raggedright
agrees
\end{minipage} \\
\midrule\noalign{}
\endhead
\bottomrule\noalign{}
\endlastfoot
12 & \texttt{Out{[}9{]}//MatrixForm=} & eta4488 & yes \\
35 & \texttt{Out{[}39{]}=} & s4by4{[}h{]}, t4by4{[}h{]}, h = 1..3 &
yes \\
74 & \texttt{Out{[}63{]}//MatrixForm=} & zeta{[}7{]} & yes \\
100 & \texttt{Out{[}81{]}//MatrixForm=} & tau{[}7{]} & yes \\
102 & \texttt{Out{[}82{]}//MatrixForm=} & Omega = sigma.tau{[}7{]} &
yes \\
106 & \texttt{Out{[}84{]}//MatrixForm=} &
tau{[}5{]}.tau{[}6{]}.tau{[}7{]} & yes \\
108 & \texttt{Out{[}85{]}//MatrixForm=} & sigma.Omega & yes \\
135 & \texttt{Out{[}98{]}//MatrixForm=} & T16\^{}A{[}8{]} & yes \\
162 & \texttt{Out{[}111{]}//MatrixForm=} & P\_L & yes \\
164 & \texttt{Out{[}112{]}//MatrixForm=} & P\_R & yes \\
313 & \texttt{Out{[}210{]}=} & T16\^{}A{[}0..7{]} & yes \\
791 & \texttt{Out{[}489{]}//MatrixForm=} &
T16\textsuperscript{A{[}1{]}\ldots T16}A{[}6{]} & yes \\
76 & \texttt{Out{[}64{]}=} & sigma ==
zeta{[}1{]}.zeta{[}2{]}.zeta{[}3{]} & yes \\
78 & \texttt{Out{[}65{]}=} & 1/2\{zeta{[}A{]}, zeta{[}B{]}\} ==
-eta{[}A,B{]} ID8, A, B = 1..7 (49 cases) & yes \\
82 & \texttt{Out{[}67{]}=} & -eta{[}A,A{]} zeta{[}A{]} ==
Transpose{[}zeta{[}A{]}{]} & yes \\
89 & \texttt{Out{[}72{]}//MatrixForm=} & sigma ==
tau{[}1{]}.tau{[}2{]}.tau{[}3{]} & yes \\
93 & \texttt{Out{[}76{]}=} & sigma.taubar{[}a{]} ==
Transpose{[}sigma.tau{[}a{]}{]} & yes \\
95 & \texttt{Out{[}77{]}//MatrixForm=} & 1/2(tau{[}A{]}.taubar{[}B{]} +
tau{[}B{]}.taubar{[}A{]}) == eta{[}A,B{]} ID8 (64 cases) & yes \\
104 & \texttt{Out{[}83{]}=} & Omega == tau{[}4{]}.tau{[}5{]}.tau{[}6{]}
& yes \\
112 & \texttt{Out{[}87{]}=} & tau{[}1{]}\ldots tau{[}7{]} == tau{[}0{]}
== ID8 & yes \\
116 & \texttt{Out{[}89{]}=} & sigma ==
tau{[}4{]}.tau{[}5{]}.tau{[}6{]}.tau{[}7{]} & yes \\
120 & \texttt{Out{[}90{]}=} &
tau{[}1{]}.tau{[}2{]}.tau{[}3{]}.tau{[}a{]} antisymmetric & yes \\
124 & \texttt{Out{[}92{]}=} &
tau{[}5{]}.tau{[}6{]}.tau{[}7{]}.taubar{[}a{]} ==
Transpose{[}tau{[}5{]}.tau{[}6{]}.tau{[}7{]}.tau{[}a{]}{]} & yes \\
126 & \texttt{Out{[}93{]}=} & tau{[}4{]}.tau{[}a{]} symmetric & yes \\
128 & \texttt{Out{[}94{]}=} & tau{[}4{]}.taubar{[}a{]} ==
Transpose{[}tau{[}4{]}.tau{[}a{]}{]} & yes \\
137 & \texttt{Out{[}99{]}=} & T0.T1.T2.T3.T{[}a{]} antisymmetric &
yes \\
139 & \texttt{Out{[}100{]}=} & T4.T5.T6.T7.T{[}a{]} symmetric, a = 0..8
& yes \\
141 & \texttt{Out{[}101{]}=} & T16\^{}A{[}8{]} == sigma16.T4.T5.T6.T7 &
yes \\
143 & \texttt{Out{[}102{]}=} & T16\^{}A{[}a{]} symmetric, a = 0..8 &
yes \\
145 & \texttt{Out{[}103{]}=} & T16\^{}A{[}a{]} antisymmetric, a = 0..8 &
yes \\
147 & \texttt{Out{[}104{]}=} & sigma16.T16\^{}A{[}a{]} antisymmetric &
yes \\
149 & \texttt{Out{[}105{]}=} & sigma16 ==
ArrayFlatten{[}\{\{-sigma,0\},\{0,sigma\}\}{]} & yes \\
157 & \texttt{Out{[}109{]}//MatrixForm=} & 1/2\{T16\^{}A{[}A{]},
T16\^{}A{[}B{]}\} == eta4488{[}{[}A+1,B+1{]}{]} ID16 (64 cases) & yes \\
159 & \texttt{Out{[}498{]}=} & T16\^{}A{[}8{]} === sigma16.T4.T5.T6.T7 &
yes \\
166 & \texttt{Out{[}113{]}=} & \{ID16 == P\_L + P\_R, P\_L.P\_L == P\_L,
P\_L.P\_R == P\_R.P\_L == ZERO16\} & yes \\
280 & \texttt{Out{[}188{]}=} & Id16 ==
T16\textsuperscript{A{[}0{]}\ldots T16}A{[}8{]} & yes \\
655 & \texttt{Out{[}392{]}=} & T16\textsuperscript{A{[}4{]}.T16}A{[}4{]}
=== -Id16 & yes \\
775 & \texttt{Out{[}476{]}=} & J = T16\^{}A{[}4{]}: \{J.J === -ID16,
sigma16 === J\^{}T.sigma16.J, !(sigma16.J antisymmetric)\} & yes \\
\end{longtable}

\hypertarget{the-same-matrices-in-every-calculation}{%
\section{The same matrices in every
calculation}\label{the-same-matrices-in-every-calculation}}

The author asked to verify that the Dirac matrices used in the
calculations are those of \texttt{4+4.nb}. Every other source of Dirac
matrices in this repository was compared with them, entry by entry:

\begin{longtable}[]{@{}
  >{\raggedright\arraybackslash}p{(\columnwidth - 2\tabcolsep) * \real{0.7500}}
  >{\raggedright\arraybackslash}p{(\columnwidth - 2\tabcolsep) * \real{0.2500}}@{}}
\toprule\noalign{}
\begin{minipage}[b]{\linewidth}\raggedright
source
\end{minipage} & \begin{minipage}[b]{\linewidth}\raggedright
equal to the \(T_a\) of 4+4.nb
\end{minipage} \\
\midrule\noalign{}
\endhead
\bottomrule\noalign{}
\endlastfoot
provenance/dirac\_matrices/author\_notebook\_T16.json (the
Pair\_Creation notebook, evaluated by WolframScript 15.0.1) & yes, all
eight \\
Revision/algebra/gammas.json (Revision record) & yes, all eight \\
universe\_model\_2/outputs/algebra.json (UM-2) & yes, all eight \\
Revision/field\_equations\_a4/wolfram/FieldEquationsA4.wl (Wolfram
package, evaluated by WolframScript) & yes, all eight \\
Revision/algebra/wolfram/RevisionAlgebra.wl (Wolfram package, evaluated
by WolframScript) & yes, all eight \\
wolfram/Dirac16ComplexAlgebra.wl (Wolfram package, evaluated by
WolframScript) & yes, all eight \\
wolfram/Dirac16ComplexGeometry.wl (Wolfram package, evaluated by
WolframScript) & yes, all eight \\
wolfram/Dirac16ComplexPrimordial.wl (Wolfram package, evaluated by
WolframScript) & yes, all eight \\
wolfram/Dirac16ComplexKohnSham.wl (Wolfram package, evaluated by
WolframScript) & yes, all eight \\
wolfram/Dirac16ComplexPairing.wl (Wolfram package, evaluated by
WolframScript) & yes, all eight \\
wolfram/Dirac16ComplexMatterAntimatter.wl (Wolfram package, evaluated by
WolframScript) & yes, all eight \\
wolfram/Dirac16Complex00.wl (Wolfram package, evaluated by
WolframScript) & yes, all eight \\
artifacts/dirac16complex/arbitrary-field/algebra-fixture.json (Stage 1
fixture) & yes, all eight \\
\end{longtable}

\texttt{Pair\_Creation\_of\_Universes\_WaveFunctionOfUniverse-4+4-Einstein-Lovelock-Nash.nb},
the author's notebook of the primordial field, builds the matrices with
the same formulas (there \(\tau\) is written directly with the formulas
that \texttt{4+4.nb} gives to its intermediate \(\zeta\);
\texttt{D2\_pair\_creation\_same\_formulas} (passed)), and its
evaluation by WolframScript in an earlier session gave
\(T_0, \dots, T_7\), \(\Gamma\), \(\sigma_{16}\), \(P_L\), \(P_R\) and
all \(S_{ab}\) equal to those of \texttt{4+4.nb}
(\texttt{D1\_pair\_creation\_extraction} (passed)). In UM-4 every
computation imports its matrices from
\texttt{universe\_model\_4/code/um4\_gammas.py}, the cell-by-cell
transcription checked here, so UM-4 uses exactly these matrices.

\hypertarget{how-to-reproduce}{%
\section{How to reproduce}\label{how-to-reproduce}}

Requirements: Python 3 with sympy
(\texttt{python\ -m\ pip\ install\ sympy}); for the PDF, pandoc and a
TeX distribution. From the repository root:

\begin{verbatim}
python universe_model_4/code/um4_dirac_matrices.py       # 55 checks; writes outputs/dirac_matrices.json and reports/dirac_matrices_checks.json
python universe_model_4/code/um4_dirac_matrices_doc.py   # writes "universe_model_4/dirac matrices.md"
python universe_model_4/tools/build_docs.py "dirac matrices"   # writes the .tex and the .pdf
\end{verbatim}

The first program needs no Mathematica: it reads the notebook file with
\texttt{universe\_model\_4/tools/nb\_reader.py}. Its run time was 38 s.

\hypertarget{appendix-the-28-generators}{%
\section{Appendix: the 28 generators}\label{appendix-the-28-generators}}

Each \(2S_{ab} = T_aT_b\) (\(a < b\)) is a signed permutation matrix; it
is listed by its non-zero entries ``row:sign column''. Rotations
(\(\eta_{aa} = \eta_{bb}\), 12 of them) satisfy
\((T_aT_b)^2 = -I_{16}\), boosts (\(\eta_{aa} = -\eta_{bb}\), 16)
\((T_aT_b)^2 = +I_{16}\).

\[
2S_{ab} = T_aT_b\qquad (0 \le a < b \le 7)
\]

\textbf{Symbols in this equation.}

\begin{itemize}
\item
  \(S_{ab}\): The scaled commutator \(S_{ab} = \tfrac14[T_a, T_b]\)
  (\([X, Y] = XY - YX\)), the generators of spin(4,4).
\item
  \(T_a\): The author's eight real \(16 \times 16\) Dirac matrices
  \texttt{T16\^{}A{[}a{]}} (\(a = 0, \dots, 7\)), the gamma matrix of
  the frame direction of \(x_a\): \(\gamma^{(x_a)} = T_a\).
\item
  \(a, b, c, d\): Frame (orthonormal) indices running over
  \(0, 1, \dots, 7\); the value \(k\) belongs to the author's coordinate
  \(x_k\) (\(x_0\) the hidden space-like direction, \(x_1, x_2, x_3\)
  ordinary space, \(x_4\) the time, \(x_5, x_6, x_7\) the extra
  time-like directions). No sum over a repeated frame index unless a sum
  is written.
\item
  \(2S_{01} = T_0T_1\) (rotation): 0:+7, 1:+6, 2:-5, 3:-4, 4:+3, 5:+2,
  6:-1, 7:-0, 8:-15, 9:-14, 10:+13, 11:+12, 12:-11, 13:-10, 14:+9,
  15:+8.
\item
  \(2S_{02} = T_0T_2\) (rotation): 0:-6, 1:+7, 2:+4, 3:-5, 4:-2, 5:+3,
  6:+0, 7:-1, 8:+14, 9:-15, 10:-12, 11:+13, 12:+10, 13:-11, 14:-8,
  15:+9.
\item
  \(2S_{03} = T_0T_3\) (rotation): 0:+5, 1:-4, 2:+7, 3:-6, 4:+1, 5:-0,
  6:+3, 7:-2, 8:-13, 9:+12, 10:-15, 11:+14, 12:-9, 13:+8, 14:-11,
  15:+10.
\item
  \(2S_{04} = T_0T_4\) (boost): 0:+5, 1:-4, 2:-7, 3:+6, 4:-1, 5:+0,
  6:+3, 7:-2, 8:-13, 9:+12, 10:+15, 11:-14, 12:+9, 13:-8, 14:-11,
  15:+10.
\item
  \(2S_{05} = T_0T_5\) (boost): 0:-6, 1:-7, 2:+4, 3:+5, 4:+2, 5:+3,
  6:-0, 7:-1, 8:+14, 9:+15, 10:-12, 11:-13, 12:-10, 13:-11, 14:+8,
  15:+9.
\item
  \(2S_{06} = T_0T_6\) (boost): 0:-7, 1:+6, 2:-5, 3:+4, 4:+3, 5:-2,
  6:+1, 7:-0, 8:+15, 9:-14, 10:+13, 11:-12, 12:-11, 13:+10, 14:-9,
  15:+8.
\item
  \(2S_{07} = T_0T_7\) (boost): 0:-0, 1:-1, 2:-2, 3:-3, 4:+4, 5:+5,
  6:+6, 7:+7, 8:+8, 9:+9, 10:+10, 11:+11, 12:-12, 13:-13, 14:-14,
  15:-15.
\item
  \(2S_{12} = T_1T_2\) (rotation): 0:+1, 1:-0, 2:+3, 3:-2, 4:+5, 5:-4,
  6:+7, 7:-6, 8:+9, 9:-8, 10:+11, 11:-10, 12:+13, 13:-12, 14:+15,
  15:-14.
\item
  \(2S_{13} = T_1T_3\) (rotation): 0:+2, 1:-3, 2:-0, 3:+1, 4:+6, 5:-7,
  6:-4, 7:+5, 8:+10, 9:-11, 10:-8, 11:+9, 12:+14, 13:-15, 14:-12,
  15:+13.
\item
  \(2S_{14} = T_1T_4\) (boost): 0:+2, 1:-3, 2:+0, 3:-1, 4:-6, 5:+7,
  6:-4, 7:+5, 8:+10, 9:-11, 10:+8, 11:-9, 12:-14, 13:+15, 14:-12,
  15:+13.
\item
  \(2S_{15} = T_1T_5\) (boost): 0:+1, 1:+0, 2:+3, 3:+2, 4:-5, 5:-4,
  6:-7, 7:-6, 8:+9, 9:+8, 10:+11, 11:+10, 12:-13, 13:-12, 14:-15,
  15:-14.
\item
  \(2S_{16} = T_1T_6\) (boost): 0:+0, 1:-1, 2:-2, 3:+3, 4:-4, 5:+5,
  6:+6, 7:-7, 8:+8, 9:-9, 10:-10, 11:+11, 12:-12, 13:+13, 14:+14,
  15:-15.
\item
  \(2S_{17} = T_1T_7\) (boost): 0:-7, 1:-6, 2:+5, 3:+4, 4:+3, 5:+2,
  6:-1, 7:-0, 8:-15, 9:-14, 10:+13, 11:+12, 12:+11, 13:+10, 14:-9,
  15:-8.
\item
  \(2S_{23} = T_2T_3\) (rotation): 0:+3, 1:+2, 2:-1, 3:-0, 4:+7, 5:+6,
  6:-5, 7:-4, 8:+11, 9:+10, 10:-9, 11:-8, 12:+15, 13:+14, 14:-13,
  15:-12.
\item
  \(2S_{24} = T_2T_4\) (boost): 0:+3, 1:+2, 2:+1, 3:+0, 4:-7, 5:-6,
  6:-5, 7:-4, 8:+11, 9:+10, 10:+9, 11:+8, 12:-15, 13:-14, 14:-13,
  15:-12.
\item
  \(2S_{25} = T_2T_5\) (boost): 0:-0, 1:+1, 2:-2, 3:+3, 4:+4, 5:-5,
  6:+6, 7:-7, 8:-8, 9:+9, 10:-10, 11:+11, 12:+12, 13:-13, 14:+14,
  15:-15.
\item
  \(2S_{26} = T_2T_6\) (boost): 0:+1, 1:+0, 2:-3, 3:-2, 4:-5, 5:-4,
  6:+7, 7:+6, 8:+9, 9:+8, 10:-11, 11:-10, 12:-13, 13:-12, 14:+15,
  15:+14.
\item
  \(2S_{27} = T_2T_7\) (boost): 0:+6, 1:-7, 2:-4, 3:+5, 4:-2, 5:+3,
  6:+0, 7:-1, 8:+14, 9:-15, 10:-12, 11:+13, 12:-10, 13:+11, 14:+8,
  15:-9.
\item
  \(2S_{34} = T_3T_4\) (boost): 0:-0, 1:-1, 2:+2, 3:+3, 4:+4, 5:+5,
  6:-6, 7:-7, 8:-8, 9:-9, 10:+10, 11:+11, 12:+12, 13:+13, 14:-14,
  15:-15.
\item
  \(2S_{35} = T_3T_5\) (boost): 0:-3, 1:+2, 2:+1, 3:-0, 4:+7, 5:-6,
  6:-5, 7:+4, 8:-11, 9:+10, 10:+9, 11:-8, 12:+15, 13:-14, 14:-13,
  15:+12.
\item
  \(2S_{36} = T_3T_6\) (boost): 0:+2, 1:+3, 2:+0, 3:+1, 4:-6, 5:-7,
  6:-4, 7:-5, 8:+10, 9:+11, 10:+8, 11:+9, 12:-14, 13:-15, 14:-12,
  15:-13.
\item
  \(2S_{37} = T_3T_7\) (boost): 0:-5, 1:+4, 2:-7, 3:+6, 4:+1, 5:-0,
  6:+3, 7:-2, 8:-13, 9:+12, 10:-15, 11:+14, 12:+9, 13:-8, 14:+11,
  15:-10.
\item
  \(2S_{45} = T_4T_5\) (rotation): 0:-3, 1:+2, 2:-1, 3:+0, 4:-7, 5:+6,
  6:-5, 7:+4, 8:-11, 9:+10, 10:-9, 11:+8, 12:-15, 13:+14, 14:-13,
  15:+12.
\item
  \(2S_{46} = T_4T_6\) (rotation): 0:+2, 1:+3, 2:-0, 3:-1, 4:+6, 5:+7,
  6:-4, 7:-5, 8:+10, 9:+11, 10:-8, 11:-9, 12:+14, 13:+15, 14:-12,
  15:-13.
\item
  \(2S_{47} = T_4T_7\) (rotation): 0:-5, 1:+4, 2:+7, 3:-6, 4:-1, 5:+0,
  6:+3, 7:-2, 8:-13, 9:+12, 10:+15, 11:-14, 12:-9, 13:+8, 14:+11,
  15:-10.
\item
  \(2S_{56} = T_5T_6\) (rotation): 0:+1, 1:-0, 2:-3, 3:+2, 4:+5, 5:-4,
  6:-7, 7:+6, 8:+9, 9:-8, 10:-11, 11:+10, 12:+13, 13:-12, 14:-15,
  15:+14.
\item
  \(2S_{57} = T_5T_7\) (rotation): 0:+6, 1:+7, 2:-4, 3:-5, 4:+2, 5:+3,
  6:-0, 7:-1, 8:+14, 9:+15, 10:-12, 11:-13, 12:+10, 13:+11, 14:-8,
  15:-9.
\item
  \(2S_{67} = T_6T_7\) (rotation): 0:+7, 1:-6, 2:+5, 3:-4, 4:+3, 5:-2,
  6:+1, 7:-0, 8:+15, 9:-14, 10:+13, 11:-12, 12:+11, 13:-10, 14:+9,
  15:-8.
\end{itemize}

\end{document}
